% !TEX program = xelatex
\documentclass[10pt, aspectratio=169]{beamer}

% --- Theme & colors (lighter Fudan-blue) ---
\usetheme{Madrid}
\usecolortheme{whale}
\setbeamercolor{structure}{fg=blue!40!white}
\setbeamercolor{palette primary}{bg=blue!30!white,fg=white}
\setbeamercolor{palette secondary}{bg=blue!25!white,fg=white}
\setbeamercolor{palette tertiary}{bg=blue!20!white,fg=white}
\setbeamercolor{palette quaternary}{bg=blue!15!white,fg=black}
\setbeamercolor{titlelike}{bg=blue!25!white,fg=black}

% --- Minimal Chinese support (only used in the title) ---
\usepackage{xeCJK}
\setCJKmainfont{FandolKai-Regular.otf}
\newcommand{\cnkai}[1]{{\CJKfontspec{FandolKai-Regular.otf}\large #1}}
\usepackage{amsmath, amssymb, amsthm, mathtools}
\usepackage{booktabs}
\usepackage{bm}
\usepackage[normalem]{ulem}
\usepackage{tikz}
\usetikzlibrary{arrows.meta, positioning, calc, fit, shapes.misc, decorations.pathreplacing}

% --- Math macros (match thesis notation; do not introduce new symbols) ---
\newcommand{\N}{\mathbb{N}}
\newcommand{\R}{\mathbb{R}}
\newcommand{\E}{\mathbb{E}}
\renewcommand{\P}{\mathbb{P}}
\newcommand{\XiLaw}{\Xi}
\newcommand{\V}{\mathcal V}
\newcommand{\M}{\mathcal M}
\newcommand{\Z}{\mathcal Z}
\newcommand{\Var}{\mathrm{Var}}
\newcommand{\e}{\mathrm{e}}
\newcommand{\one}{\mathbf{1}}

% --- Use plain numeric labels (not ball-enumitem circles) ---
\setbeamertemplate{enumerate items}[default]
\setbeamertemplate{section in toc}{\inserttocsectionnumber.~\inserttocsection\par}

% Thesis environment used by some snippets.
% beamer defines theorem/lemma/corollary already; override Chinese names from ctex.
\newtheorem{assumption}{Assumption}
\newtheorem{proposition}[theorem]{Proposition}

% --- Compact slide style ---
\setbeamersize{text margin left=8mm, text margin right=8mm}
\setlength{\parskip}{3pt}

% --- Title page metadata ---
\title[Asymptotic speed of $N$-BRW]{Asymptotic speed of the $N$-branching random walk\\with infinite offspring\\[6pt] \cnkai{可无限分支的 $N$-分支随机游走的渐近速率性质}\\[6pt] \small Master's thesis defense}
\author[Wang Yiyang]{Presenter: Yiyang Wang \\ \small Instructor: Prof.~Élie Aïdékon}
\institute[Fudan University]{School of Mathematical Sciences, Fudan University}
\date{May 20, 2026}

\begin{document}

% --- Title ---
\begin{frame}
    \titlepage
\end{frame}

% --- Outline ---
\begin{frame}{Outline}
    \tableofcontents
\end{frame}

% ==========================================
\section{Background and model}
% ==========================================

\begin{frame}{BRW on the marked tree space}
    % Purpose: formalize BRW as random marked tree, selection as deterministic functional.
    \small
    % Ulam--Harris trees, marked trees, and indexed marked tree families.
\begin{itemize}
    \item \textbf{Ulam--Harris vertex space:}
    $\mathbb V:=\bigcup_{n=0}^\infty\mathbb N^n$, $\mathbb N^0:=\{\emptyset\}$ (root).
    $|u|$ = generation of $u$.
    \item \textbf{Rooted labelled tree:} $T\subset\mathbb V$ s.t.
    (i) $\emptyset\in T$;
    (ii) ancestors closed;
    (iii) $uj\in T\Rightarrow uj'\in T$ for $j'\le j$.
    \\ 
    Space: $\mathbb T$; $\sigma$-algebra $\mathcal G_{\mathbb T}$ from cylinders $\{u\in T\}$ (coarser than Borel $\sigma$-algebra).
    \item \textbf{Marked tree:} $(T,V)$, $V:T\to\R$. \\
    Space: $\mathbb M:=\{(T,V):T\in\mathbb T,\,V:T\to\R\}$;
    $\mathcal G_{\mathbb M}$ from $\{u\in T,\,V(u)\in A\}$, $A\in\mathcal B(\R)$.
    \item \textbf{BRW:} $(\mathcal T,\mathcal V)\in(\mathbb M,\mathcal G_{\mathbb M})$, denoted $\mathcal V$.
    \item \textbf{Indexed marked tree family:} $\boldsymbol T\in\mathbb T^{\mathbb I}$, $V:\{(u,i):u\in\boldsymbol T(i),\,i\in\mathbb I\}\to\R$.\\
    Space: $\mathbb M^{\mathbb I}$; $\mathcal G_{\mathbb M}^{\otimes\mathbb I}$ from $\{u\in\boldsymbol T(i),\,V(u,i)\in A\}$.
    \item \textbf{Selection operator:} $\Phi_N:\mathbb M^{\mathbb I}\to\mathbb M^{\mathbb I}$, keeps $N$ leftmost per generation.\\
    $N$-BRW: $\mathcal V^N:=\Phi_N(\mathcal V)$.
    \item \textbf{Position \& speed as functionals:}
    $\mathcal M_n$ = leftmost; $\mathcal M_n^{N|a}$ for $N$-BRW.\\
    $v_N:=\lim\frac{1}{n}\mathcal M_n^{N|a}$, $v_\infty:=\lim\frac{1}{n}\mathcal M_n$.
\end{itemize}

\end{frame}

\begin{frame}{What is $N$-BRW?}
    % Purpose: N-BRW definition with illustration, selection mechanism, and notation.
    \begin{columns}[T]
        \column{0.46\textwidth}
            \centering
            \begin{tikzpicture}[
    xscale=1.2,
    yscale=0.55,
    survivor/.style={circle, fill=blue!70!black, inner sep=1.2pt},
    killed/.style={circle, draw=black, fill=white, inner sep=1.2pt},
    live-edge/.style={thick, blue!70!black},
    killed-edge/.style={dashed, black},
]
    \draw[->, >=Stealth, black, thick] (0, -3) -- (0, 2.5);
    \draw[-, black, thick, dashed] (2, -3) -- (2, 2.5);
    \draw[-, black, thick, dashed] (4, -3) -- (4, 2.5);

    % --- Generation 0 ---
    \node[survivor] (root) at (0, 0) {};
    \node[anchor=south] at (0, 3.0) {$n=0$};

    % --- Generation 1 (4 offspring per parent) ---
    \node[survivor] (g1-1) at (2, -1.6) {};
    \node[survivor] (g1-2) at (2, -0.3) {};
    \node[survivor] (g1-3) at (2, 0.5)  {};
    \node[killed]   (g1-4) at (2, 1.8)  {};

    \draw[live-edge] (root) -- (g1-1);
    \draw[live-edge] (root) -- (g1-2);
    \draw[live-edge] (root) -- (g1-3);
    \draw[killed-edge] (root) -- (g1-4);

    \node[anchor=south] at (2, 3.0) {$n=1$};

    % --- Generation 2 (each survivor produces 4 offspring) ---
    \node[survivor] (g2-1-1) at (4, -2.6) {};
    \node[survivor] (g2-1-2) at (4, -1.5) {};
    \node[killed]   (g2-1-3) at (4, 0.4) {};
    \node[killed]   (g2-1-4) at (4, 0.9) {};
    \foreach \i in {1,2} \draw[live-edge] (g1-1) -- (g2-1-\i);
    \foreach \i in {3,4} \draw[killed-edge] (g1-1) -- (g2-1-\i);

    \node[survivor] (g2-2-1) at (4, -1.8) {};
    \node[killed]   (g2-2-2) at (4, -0.4) {};
    \node[killed]   (g2-2-3) at (4, -0.8) {};
    \node[killed]   (g2-2-4) at (4, 0.4) {};
    \draw[live-edge] (g1-2) -- (g2-2-1);
    \foreach \i in {2,3,4} \draw[killed-edge] (g1-2) -- (g2-2-\i);

    \foreach \y/\i in {-0.6/1, -0.3/2, 1.2/3, 1.9/4} {
        \node[killed] (g2-3-\i) at (4, \y) {};
        \draw[killed-edge] (g1-3) -- (g2-3-\i);
    }

    \node[anchor=south] at (4, 3.0) {$n=2$};
\end{tikzpicture}

            \vspace{1mm}

            \footnotesize A 3-BRW: each generation keeps the $N=3$ leftmost particles.
        \column{0.54\textwidth}
            \small
            % $N$-BRW: selection mechanism and key notation.
\begin{itemize}
    \item \textbf{Init:} $\lfloor N^a\rfloor$ particles at $0$ ($a\in[0,1]$).
    \item \textbf{Branch:} each reproduces via displacement p.p.\ $\Xi$.
    \item \textbf{Selection} (keep $N$ leftmost): a deterministic functional
    $\Phi_N:\mathbb M^{\mathbb I}\to\mathbb M^{\mathbb I}$;
    $\mathcal V^N:=\Phi_N(\mathcal V)$ ($N$-BRW).
    \item \textbf{Leftmost position:} $\mathcal M_n^{N|a}$ at generation $n$.
    \item \textbf{Speed:} $v_N:=\lim_{n\to\infty}\frac{1}{n}\mathcal M_n^{N|a}$
    (exists via Kingman); $v_\infty$ = speed of unconstrained BRW.
\end{itemize}

    \end{columns}
\end{frame}

\begin{frame}{Motivation: Brunet--Derrida speed scaling}
    % Purpose: introduce topic + the conjectural scaling formula.
    \begin{block}{Brunet--Derrida conjecture (1997) and refined conjecture (2007)}
        % Extracted verbatim (English part) from: docs/main.tex (Future Work item 1; BD refined conjecture form)
% Purpose: state the conjectured speed correction scaling law.
\begin{equation*}
    v_\infty - v_N
    \asymp
    \frac{1}{\log^2 N},
    \qquad
    \text{and conjecturally}\quad
    v_\infty - v_N
    =
    C{\left(\log N+3\log\log N+o_N(\log\log N)\right)^{-2}}.
\end{equation*}

    \end{block}
    \vspace{-1mm}
    \begin{itemize}
        \item Goal here: a rigorous $(\log N)^{-2}$ law with an explicit leading constant.
        \item Thesis contribution: extends the law to general branching laws $\Xi$ (possibly infinite offspring), for the $\alpha=2$ (Gaussian tail) case.
    \end{itemize}
\end{frame}

% ==========================================
\section{Main theorem: asymptotic speed}
% ==========================================

\begin{frame}{Main theorem (asymptotic speed)}
    % Purpose: state the main theorem from the thesis (no new notation).
    % Theorem `theorem:speed` (English-only, faithful to contents/part-5-eng.tex).
\begin{theorem}[Asymptotic speed of the leftmost particle]
Consider the $N$-branching random walk $\mathcal V^{N|a}$ that starts with $\lfloor N^a\rfloor$ particles at position $0$, with $a\in[0,1]$, and whose branching law $\Xi$ satisfies the standing assumptions in the thesis (supercriticality, boundary case, finite variance, finite cross second moment, and moment conditions on $\Xi_1$).
Then the speed of the leftmost particle exists and
\[
    v_N
    =
    \lim_{n\to\infty}\frac{1}{n}\mathcal M^{N|a}_n
    =
    \frac{\pi^2}{2}\sigma^2\frac{1}{\log^2 N}
    +
    o_N\!\left(\frac{1}{\log^2 N}\right).
\]
\end{theorem}

\end{frame}

\begin{frame}{Assumptions}
    % Purpose: list the assumptions used in the main theorem, with full moment details.
    \small
    \textbf{Log-Laplace transform:}
    \[
        \psi(\lambda):=\log\E\Big[\sum_{x\in\Xi}\e^{-\lambda x}\Big],\qquad \lambda\in\R.
    \]
    \vspace{-2mm}
    \begin{enumerate}
        \item \textbf{Supercriticality \& viability:} $\psi(0)>0$ (supercritical) and $\#\Xi\ge 1$ a.s.\ (viability: at least one child).
        \item \textbf{Boundary case:} $\psi(1)=0$, $\psi'(1)=0$ (after centering, the spine walk is centered).
        \item \textbf{Finite variance:} $\sigma^2:=\psi''(1)\in(0,\infty)$.
        \item \textbf{Moment control on the leftmost step $\Xi_1$:}
        \begin{itemize}
            \item $\E[(\Xi_1)_+] < \infty$ \quad (finite positive part),
            \item {\color{gray}\sout{$\E[(\Xi_1)_+^4] < \infty$}} \quad {\color{gray} (removed --- no longer needed after refinement)},
            \item {\color{gray!60!white} $\exists\,\lambda>0,\; \E[\e^{\lambda\Xi_1}] < \infty$ \quad (finite exponential moment --- only for Theorem~1.1)}.
        \end{itemize}
        \item \textbf{Cross second moment (Assumption~6):}
        \[
            \log \E\!\left[\sum_{i\neq j}\e^{-(\Xi_i+\Xi_j)}\right] < \infty,
        \]
        controls the second moment of tube-counting functionals, used in Lemma~3.8 and the Paley--Zygmund step.
    \end{enumerate}
\end{frame}

\begin{frame}{What is the \emph{boundary case}? (meaning and key properties)}
    \small
    \begin{itemize}
        \item The \textbf{boundary case} is the critical tilt where
        \[
            \psi(1)=0
            \quad\text{and}\quad
            \psi'(1)=0,
        \]
        so the size-biased (spine) walk is \textbf{centered}.
        \item \textbf{Leftmost position diverges:}
        \[
            \lim_{n\to\infty}\mathcal M_n = +\infty
            \quad\text{a.s.}
        \]
        \item \textbf{Vanishing speed of the unconstrained BRW:} $v_\infty = \lim_{n\to\infty}\frac{1}{n}\mathcal M_n = 0$.
        \item The (critical) \textbf{additive martingale} $W_n:=\sum_{|x|=n}\e^{-\mathcal V(x)}$ has $\E[W_n]=1$,
        but in the boundary case it typically satisfies $W_n\to 0$ a.s.\ under mild conditions.
        \item A non-degenerate renormalization is given by the \textbf{derivative martingale}
        $D_n:=\sum_{|x|=n}\mathcal V(x)\e^{-\mathcal V(x)}$ (central object for extremal positions).
        \item \textbf{Generality:} in most cases a BRW can be transformed to satisfy the boundary case (Jaffuel, 2009).
    \end{itemize}
\end{frame}

\begin{frame}{Existence of $v_N$}
    % Purpose: state the existence result (Proposition 2 + Corollary) from the thesis.
    \small
    \begin{block}{Proposition}
        Under $\psi(1)=0$ and $\E[(\Xi_1)_+]<\infty$, there exists $v_N=v_N(\Xi)\in\R$ such that,
        for the $N$-BRW $\mathcal V^N$ with $\mathcal Z^N_0=N\delta_0$,
        \[
            \lim_{n\to\infty}\frac{1}{n}\mathcal M^N_n = v_N
            \qquad\text{a.s.\ and in }\mathrm L^1(\P).
        \]
    \end{block}
    {\small Proved via Kingman's subadditive ergodic theorem.}
    \vspace{1mm}
    \begin{corollary}
        If additionally $\psi(0)>0$ (supercritical), the same holds for any $\mathcal Z^N_0\in\mathcal C(N_0)$ with $N_0\le N$.
    \end{corollary}
\end{frame}

% ==========================================
\section{Upper bound: restart stage and proof}
% ==========================================

\begin{frame}{Upper comparison construction}
    \begin{columns}
        \column{0.35\textwidth}
        \small
        Kill all but one at $0$; evolve as $N$-BRW for $\ell_N$; keep only leftmost at $k\ell_N$, repeat.
        \column{0.65\textwidth}
        \centering
        \resizebox{\textwidth}{!}{% Source: contents/part-5-eng.tex (Figure: construction of V^{upper})
\begin{tikzpicture}[
    xscale=2,
    yscale=1.2,
    >=Stealth,
    % styles
    survivor/.style={circle, fill=blue!70!black, inner sep=1.2pt},
    killed-node/.style={circle, draw=blue!70!black, fill=white, inner sep=1.2pt, cross out, thick, minimum size=4pt},
    edge-live/.style={thick, blue!70!black},
    edge-killed/.style={thick, dashed, blue!30},
    cloud-tri/.style={fill=blue!10, opacity=0.5},
    ref-line/.style={thin, gray!40, dashed}
]

    % time reference lines (0, l_N, 2l_N, 3l_N)
    \foreach \x/\label in {0/0, 2/\ell_N, 4/2\ell_N, 6/3\ell_N} {
        \draw[ref-line] (\x, -0.5) -- (\x, 4.5);
        \node[below] at (\x, -0.5) {$\label$};
    }

    % ===== Block 1: 0 -> l_N
    \node[survivor] (n0) at (0, 0) {};
    \node[left] at (n0) {$\emptyset$};

    \fill[cloud-tri] (0, 0) -- (2, 0.8) -- (2, 2.2) -- cycle;
    \draw[edge-live] (0,0) -- (0.8, 0.3) -- (1.4, 0.6) -- (2, 0.8);
    \draw[edge-killed] (0.8, 0.3) -- (1.3, 1.2) -- (2, 1.3);
    \draw[edge-killed] (1.4, 0.6) -- (1.8, 1.5) -- (2, 1.8);
    \draw[edge-killed] (0,0) -- (0.9, 1.2) -- (2, 2.2);

    \node[survivor] (m1) at (2, 0.8) {};
    \node[killed-node] at (2, 1.3) {};
    \node[killed-node] at (2, 1.8) {};
    \node[killed-node] at (2, 2.2) {};

    % ===== Block 2: l_N -> 2l_N
    \fill[cloud-tri] (2, 0.8) -- (4, 1.8) -- (4, 3.2) -- cycle;

    \draw[edge-live] (2, 0.8) -- (2.5, 1.1) -- (3.0, 1.4) -- (3.5, 1.3) -- (4, 1.8);
    \draw[edge-killed] (2.5, 1.1) -- (3.2, 2.0) -- (4, 2.3);
    \draw[edge-killed] (3.0, 1.4) -- (3.6, 2.1) -- (4, 2.8);
    \draw[edge-killed] (2, 0.8) -- (2.8, 2.2) -- (4, 3.2);

    \node[survivor] (m2) at (4, 1.8) {};
    \node[killed-node] at (4, 2.3) {};
    \node[killed-node] at (4, 2.8) {};
    \node[killed-node] at (4, 3.2) {};

    % ===== Block 3: 2l_N -> 3l_N
    \fill[cloud-tri] (4, 1.8) -- (6, 2.7) -- (6, 4.1) -- cycle;

    \draw[edge-live] (4, 1.8) -- (4.5, 2.0) -- (5.0, 2.4) -- (5.5, 2.3) -- (6, 2.7);
    \draw[edge-killed] (4.5, 2.0) -- (5.2, 2.9) -- (6, 3.2);
    \draw[edge-killed] (5.0, 2.4) -- (5.6, 3.1) -- (6, 3.7);
    \draw[edge-killed] (4, 1.8) -- (4.8, 3.2) -- (6, 4.1);

    \node[survivor] (m3) at (6, 2.7) {};
    \node[killed-node] at (6, 3.2) {};
    \node[killed-node] at (6, 3.7) {};
    \node[killed-node] at (6, 4.1) {};

    % labels
    \node[below, blue!70!black] at (2, 0.8) {$\mathcal{M}^N_{\ell_N}$};
    \node[below, blue!70!black] at (4, 1.8) {$\mathcal{M}^N_{2\ell_N}$};
    \node[below, blue!70!black] at (6, 2.7) {$\mathcal{M}^N_{3\ell_N}$};
\end{tikzpicture}
}
    \end{columns}
\end{frame}

\begin{frame}{Coupling / monotonicity}
    \small
    The construction yields $\mathcal M_{k\ell_N}^{N|a}\le \mathcal M_{k\ell_N}^{\mathrm{upper}}$ for all $k$.
    \vspace{1mm}
    
    \textbf{Partial order $\succcurlyeq$:}
    $\mu_1\succcurlyeq\mu_2$ iff $M(\mu_1)\le M(\mu_2)$ and $x_i\ge y_i$ for $i=1,\dots,M(\mu_1)$ (fewer particles, shifted right).
    \vspace{1mm}
    
    % Extracted from: contents/part-4-eng.tex (Corollary \label{corollary:2})
\begin{corollary}[Coupling / monotonicity]
For $N_1\le N_2$ and $\mu_1\succcurlyeq\mu_2$ with $\mu_1\in\mathcal C(N_1)$, $\mu_2\in\mathcal C(N_2)$, there exists a coupling $(Z^1_n,Z^2_n)_{n\ge 0}$ of $\mathcal V^{N_1}$ and $\mathcal V^{N_2}$ such that $Z^1_0=\mu_1$, $Z^2_0=\mu_2$ and $Z^1_n\succcurlyeq Z^2_n$ a.s.\ for all $n\ge 0$. In particular, $\min Z^1_n\ge\min Z^2_n$ a.s.
\end{corollary}

    \vspace{1mm}
    \centering
    \resizebox{0.58\textwidth}{!}{% Source: contents/part-4-eng.tex (coupling illustration: 3-BRW vs dominating 2-BRW)
\begin{tikzpicture}[
    xscale=1,
    yscale=1.5,
    >=Stealth,
    % styles
    node-both/.style={circle, fill=blue!70!black, draw=green!60!black, line width=1.2pt, inner sep=1.2pt},
    node-N/.style={circle, fill=blue!70!black, inner sep=1.8pt},
    node-k/.style={circle, fill=green!60!black, inner sep=1.8pt},
    node-N-dead/.style={circle, draw=blue!70!black, fill=white, inner sep=1.4pt, cross out, draw=blue!80!black, line width=1.2pt},
    node-k-dead/.style={circle, draw=green!60!black, fill=white, inner sep=1.4pt, cross out, draw=green!60!black, line width=1.2pt},
    node-dead/.style={circle, draw=black!40, fill=white, inner sep=1.4pt},
    edge-N/.style={thick, blue!70!black},
    edge-k/.style={thick, green!60!black},
    edge-N-dashed/.style={dashed, blue!70!black},
    edge-k-dashed/.style={dashed, green!60!black},
    edge-both/.style={
        thick,
        blue!70!black,
        postaction={
            draw,
            green!60!black,
            dash pattern=on 5pt off 5pt
        }
    },
    edge-dashed/.style={dashed, black!30},
    kill-line/.style={red!80!black, dotted, ultra thick}
]

    \foreach \y in {0,1} \draw[-, black!20, dashed] (-4.5, \y) -- (2.5, \y);
    \node[right=10pt] at (2.7, 0) {$n=0$};
    \node[right=10pt] at (2.7, 1) {$n=1$};

    % n=0
    \node[node-k] (n0-k-1) at (-1.6, 0) {};
    \node[below] at (n0-k-1) {{\color{green!60!black}$\Xi_1$}};
    \node[node-k] (n0-k-2) at (2.2, 0) {};
    \node[below] at (n0-k-2) {{\color{green!60!black}$\Xi_2$}};
    \node[node-N] (n0-N-1) at (-2.2, 0) {};
    \node[below] at (n0-N-1) {{\color{blue!70!black}$\Xi_1$}};
    \node[node-N] (n0-N-2) at (-0.8, 0) {};
    \node[below] at (n0-N-2) {{\color{blue!70!black}$\Xi_2$}};
    \node[node-N] (n0-N-3) at (-0.2, 0) {};
    \node[below] at (n0-N-3) {{\color{blue!70!black}$\Xi_3$}};

    % n=1
    \node[node-k] (n1-k-11) at (-3.0, 1) {};
    \draw[edge-k] (n0-k-1) -- (n1-k-11);
    \node[node-k] (n1-k-12) at (-2.1, 1) {};
    \draw[edge-k] (n0-k-1) -- (n1-k-12);
    \node[node-k-dead] (n1-k-13) at (-1.3, 1) {};
    \draw[edge-k-dashed] (n0-k-1) -- (n1-k-13);

    \node[node-N] (n1-N-11) at (-3.6, 1) {};
    \draw[edge-N] (n0-N-1) -- (n1-N-11);
    \node[node-N] (n1-N-12) at (-2.7, 1) {};
    \draw[edge-N] (n0-N-1) -- (n1-N-12);
    \node[node-N-dead] (n1-N-13) at (-1.9, 1) {};
    \draw[edge-N-dashed] (n0-N-1) -- (n1-N-13);

    \node[node-N] (n1-N-21) at (-4.0, 1) {};
    \draw[edge-N] (n0-N-2) -- (n1-N-21);
    \node[node-k-dead] (n1-k-21) at (-1.0, 1) {};
    \draw[edge-k-dashed] (n0-k-2) -- (n1-k-21);

    \node[node-N-dead] (n1-N-31) at (-0.4, 1) {};
    \draw[edge-N-dashed] (n0-N-3) -- (n1-N-31);
    \node[node-N-dead] (n1-N-32) at (0.8, 1) {};
    \draw[edge-N-dashed] (n0-N-3) -- (n1-N-32);
    \node[node-N-dead] (n1-N-33) at (1.9, 1) {};
    \draw[edge-N-dashed] (n0-N-3) -- (n1-N-33);
\end{tikzpicture}
}
\end{frame}

\begin{frame}{Upper bound: SLLN}
    % Purpose: derivation leading to v_N bound, then SLLN verification.
    \small
    % Extracted verbatim (English part) from: contents/part-5-eng.tex (upper bound inequality chain)
% Purpose: derive v_N upper bound from V^{upper} + SLLN + Proposition 3.1.
\begin{align*}
    v_N
    &=
    \lim_{n\to\infty}
    \frac{1}{n}\mathcal M^{N|a}_n
    =
    \lim_{k\to\infty}
    \frac 1{k\ell_N}
    \mathcal M^{N|a}_{k\ell_N}
    \\ &
    \le
    \limsup_{k\to\infty}
    \frac 1{k\ell_N}
    \mathcal M^{\mathrm{upper}}_{k\ell_N}
    =
    \frac 1{\ell_N}
    \limsup_{k\to\infty}
    \frac 1{k}
    \left(
        \sum_{i=1}^k
        \left(
            \mathcal M^{\mathrm{upper}}_{i\ell_N}
            -
            \mathcal M^{\mathrm{upper}}_{(i-1)\ell_N}
        \right)
    \right)
    \\ &
    =
    \frac 1{\ell_N}
    \limsup_{k\to\infty}
    \frac 1{k}
    \left(
        \sum_{i=1}^k
        \min
        \left(
            \mathcal V^{\mathrm{upper}}_{i\ell_N}
            -
            \mathcal M^{\mathrm{upper}}_{(i-1)\ell_N}
        \right)
    \right)
    \\ &
    =
    \frac 1{\ell_N}
    \mathbb E\left[
        \min\mathcal V^{N|0}_{\ell_N}
    \right]
    =
    \frac 1{\ell_N}
    \mathbb E\left[
        \mathcal M^{N|0}_{\ell_N}
    \right].
\end{align*}

    The increments $\min(\mathcal V^{\mathrm{upper}}_{(k+1)\ell_N}-\mathcal M^{\mathrm{upper}}_{k\ell_N})$ are i.i.d.~$\sim\min\mathcal V^{N|0}_{\ell_N}$; SLLN applies since
    $\E[\mathcal M^{N|0}_{\ell_N}]<\infty$ (by \textbf{Assumption}: $\E[(\Xi_1)_+]<\infty$).
\end{frame}

\begin{frame}{Upper bound: what needs to be proved}
    % Purpose: reduce to E[M_{ℓ_N}], state Theorem 1.1 with conditions, motivate restart.
    \small
    The chain gives $v_N\le\frac{1}{\ell_N}\E[\mathcal M_{\ell_N}^{N|0}]$.
    Fix $t>0$, choose $\ell_N=\ell_N(t)$ with $\ell_N=t\log^3 N+o(\log^3 N)$.
    Since we may let $t\to\infty$ eventually, it suffices to show
    \[
        \E[\mathcal M_{\ell_N}^{N|0}]
        \le
        \frac{\pi^2}{2}\sigma^2\,t\log N
        +
        O_N(\log N),
    \]
    which yields $v_N\le\frac{\pi^2}{2}\sigma^2\frac{1}{\log^2 N}+o_N\!\left(\frac{1}{\log^2 N}\right)$ after $t\to\infty$.

    \medskip
    \begin{theorem}[{Theorem~1.1, $\mathrm{L}^2$ edition, $a\in(0,1]$}]
        Fix $a\in(0,1]$ and $t>0$, let $\ell_N$ satisfy $\ell_N/\log^3 N\to t$. Under Assumptions~1--8 (supercritical, viability, boundary case, $\sigma^2<\infty$, $\E[(\Xi_1)_+]<\infty$, $\E[(\Xi_1)_+^4]<\infty$, cross second moment),
        \[
            \frac{1}{\log N}\min_{|u|=\ell_N}\mathcal V^{N|a}(u)
            \;\xrightarrow{\mathrm{L}^2}\;
            -a+\frac{\pi^2}{2}\sigma^2 t.
        \]
    \end{theorem}

    \textbf{But:} for $a=0$, Theorem~1.1 requires an extra exponential moment $\E[\e^{\lambda\Xi_1}]<\infty$ (\textbf{undesirable}).
    The restart mechanism (next slide) removes it: a ``backup'' particle supplies $\lfloor N^\alpha\rfloor$ offspring ($\alpha>0$), reducing $a=0$ to the $a=\alpha>0$ case.
\end{frame}

\begin{frame}{Restart construction (1): operator \& result}
    % Purpose: Part 1 — restart operator, \kappa, \tau, K, T.
    \small
    % Part 1 of restart construction definitions: truncated branching & split time.
% The restart mechanism \Phi_{\mathrm{re}} is a deterministic operator on the marked tree space.
% Composing \Phi_{\mathrm{re}} with the original BRW yields the coupled killed BRW \mathcal V^{\mathrm b}.
\centering
\resizebox{0.80\textwidth}{!}{% Source: contents/part-5-eng.tex (Figure: coupling restart mechanism / restart illustration)
\begin{tikzpicture}[
            xscale=2, 
            yscale=1.2, 
            >=Stealth,
            % 样式定义
            survivor/.style={circle, fill=blue!70!black, inner sep=1.2pt},
            killed-node/.style={circle, draw=blue!70!black, fill=white, inner sep=1.2pt, cross out, thick, minimum size=4pt},
            edge-live/.style={thick, blue!70!black},
            edge-killed/.style={thick, dashed, blue!30},
            cloud-tri/.style={fill=blue!10, opacity=0.5},
            ref-line/.style={thin, gray!40, dashed}
        ]

            % 1. 时间周期参考线 (0, l_N, 2l_N, 3l_N)
            \foreach \x/\label in {0/0, 2/\ell_N, 4/2\ell_N, 6/3\ell_N} {
                \draw[ref-line] (\x, -0.5) -- (\x, 4.5);
                \node[below] at (\x, -0.5) {$\label$};
            }

            % ================= 第一周期: 0 到 l_N =================
            % 起点
            \node[survivor] (n0) at (0, 0) {};
            \node[left] at (n0) {$\emptyset$};

            % 三角形阴影 (顶点在 n0)
            % 假设在 l_N 时刻，粒子分布在 [0.8, 2.2] 之间
            \fill[cloud-tri] (0,0) -- (2, 0.8) -- (2, 2.2) -- cycle;

            % 高频分支
            \draw[edge-live] (0,0) -- (0.5, 0.2) -- (1.0, 0.5) -- (1.5, 0.4) -- (2, 0.8);
            \draw[edge-killed] (0.5, 0.2) -- (1.2, 1.2) -- (2, 1.3);
            \draw[edge-killed] (1.0, 0.5) -- (1.6, 1.0) -- (2, 1.8);
            \draw[edge-killed] (0,0) -- (0.8, 0.9) -- (2, 2.2);

            % 瓶颈处粒子
            \node[survivor] (m1) at (2, 0.8) {}; 
            \node[killed-node] at (2, 1.3) {};
            \node[killed-node] at (2, 1.8) {};
            \node[killed-node] at (2, 2.2) {};

            % ================= 第二周期: l_N 到 2l_N =================
            % 起点是 m1 (2, 0.8)
            % 假设在 2l_N 时刻，粒子分布在 [1.8, 3.2] 之间
            \fill[cloud-tri] (2, 0.8) -- (4, 1.8) -- (4, 3.2) -- cycle;

            % 高频分支
            \draw[edge-live] (2, 0.8) -- (2.5, 1.1) -- (3.0, 1.4) -- (3.5, 1.3) -- (4, 1.8);
            \draw[edge-killed] (2.5, 1.1) -- (3.2, 2.0) -- (4, 2.3);
            \draw[edge-killed] (3.0, 1.4) -- (3.6, 2.1) -- (4, 2.8);
            \draw[edge-killed] (2, 0.8) -- (2.8, 2.2) -- (4, 3.2);

            % 瓶颈处粒子
            \node[survivor] (m2) at (4, 1.8) {};
            \node[killed-node] at (4, 2.3) {};
            \node[killed-node] at (4, 2.8) {};
            \node[killed-node] at (4, 3.2) {};

            % ================= 第三周期: 2l_N 到 3l_N =================
            % 起点是 m2 (4, 1.8)
            % 假设在 3l_N 时刻，粒子分布在 [2.7, 4.1] 之间
            \fill[cloud-tri] (4, 1.8) -- (6, 2.7) -- (6, 4.1) -- cycle;

            % 高频分支
            \draw[edge-live] (4, 1.8) -- (4.5, 2.0) -- (5.0, 2.4) -- (5.5, 2.3) -- (6, 2.7);
            \draw[edge-killed] (4.5, 2.0) -- (5.2, 2.9) -- (6, 3.2);
            \draw[edge-killed] (5.0, 2.4) -- (5.6, 3.1) -- (6, 3.7);
            \draw[edge-killed] (4, 1.8) -- (4.8, 3.2) -- (6, 4.1);

            % 瓶颈处粒子
            \node[survivor] (m3) at (6, 2.7) {};
            \node[killed-node] at (6, 3.2) {};
            \node[killed-node] at (6, 3.7) {};
            \node[killed-node] at (6, 4.1) {};

            % 瓶颈标注
            \node[above right] at (m1) {\scriptsize keep leftmost};
            \node[above right] at (m2) {\scriptsize keep leftmost};
            \node[above right] at (m3) {\scriptsize keep leftmost};
\end{tikzpicture}

}

\vspace{4pt}
\begin{itemize}
    \item \textbf{Restart operator} $\Phi_{\mathrm{re}}:\mathbb M\to\mathbb M$:
    a deterministic mapping on marked trees; composing it with the original BRW $\mathcal V$ yields
    the coupled killed BRW $\mathcal V^{\mathrm b}$.
    \item $M>0$ (indep.~of $N$) s.t.\ $p_0:=\mathbb P(\#\Xi^{(M)\mathrm b}\ge 2)>0$,
    $m:=\mathbb E[\#\Xi^{(M)}]>1$, with
    $\Xi^{(M)\mathrm b}:=\delta_{\Xi_1}{\bf 1}_{\{\#\Xi\ge 1\}}+\delta_{\Xi_2}{\color{red}{\bf 1}_{\{\Xi_2\le M\}}}{\bf 1}_{\{\#\Xi\ge 2\}}$,
    $\Xi^{(M)}:=\delta_{\Xi_1}{\color{red}{\bf 1}_{\{\Xi_1\le M\}}}{\bf 1}_{\{\#\Xi\ge 1\}}+\delta_{\Xi_2}{\color{red}{\bf 1}_{\{\Xi_2\le M\}}}{\bf 1}_{\{\#\Xi\ge 2\}}$.
    \item $\tau_1:=\inf\{n\ge 0:\ \#\mathcal Z^{\mathrm b}_{n+1}\ge 2\}$.
    Lineage: $w_n$ ($0\le n\le\tau_1$), $\rho_1:=w_{\tau_1}$.
    \item $v_1$: left-most child of $\rho_1$; $u_1$: second left-most.
\end{itemize}

\end{frame}

\begin{frame}{Restart construction (2): truncated branching \& first split}
    % Purpose: Part 2 — M, \Xi^{(M)b}, \Xi^{(M)}, \tau_1, \rho_1, v_1, u_1, post-split evolution.
    \small
    % Part 2 of restart construction definitions: iteration, \kappa, \tau, K, T.
\centering
\resizebox{0.80\textwidth}{!}{% Source: contents/part-5-eng.tex (Figure: coupling restart mechanism / restart illustration)
\begin{tikzpicture}[
            xscale=2, 
            yscale=1.2, 
            >=Stealth,
            % 样式定义
            survivor/.style={circle, fill=blue!70!black, inner sep=1.2pt},
            killed-node/.style={circle, draw=blue!70!black, fill=white, inner sep=1.2pt, cross out, thick, minimum size=4pt},
            edge-live/.style={thick, blue!70!black},
            edge-killed/.style={thick, dashed, blue!30},
            cloud-tri/.style={fill=blue!10, opacity=0.5},
            ref-line/.style={thin, gray!40, dashed}
        ]

            % 1. 时间周期参考线 (0, l_N, 2l_N, 3l_N)
            \foreach \x/\label in {0/0, 2/\ell_N, 4/2\ell_N, 6/3\ell_N} {
                \draw[ref-line] (\x, -0.5) -- (\x, 4.5);
                \node[below] at (\x, -0.5) {$\label$};
            }

            % ================= 第一周期: 0 到 l_N =================
            % 起点
            \node[survivor] (n0) at (0, 0) {};
            \node[left] at (n0) {$\emptyset$};

            % 三角形阴影 (顶点在 n0)
            % 假设在 l_N 时刻，粒子分布在 [0.8, 2.2] 之间
            \fill[cloud-tri] (0,0) -- (2, 0.8) -- (2, 2.2) -- cycle;

            % 高频分支
            \draw[edge-live] (0,0) -- (0.5, 0.2) -- (1.0, 0.5) -- (1.5, 0.4) -- (2, 0.8);
            \draw[edge-killed] (0.5, 0.2) -- (1.2, 1.2) -- (2, 1.3);
            \draw[edge-killed] (1.0, 0.5) -- (1.6, 1.0) -- (2, 1.8);
            \draw[edge-killed] (0,0) -- (0.8, 0.9) -- (2, 2.2);

            % 瓶颈处粒子
            \node[survivor] (m1) at (2, 0.8) {}; 
            \node[killed-node] at (2, 1.3) {};
            \node[killed-node] at (2, 1.8) {};
            \node[killed-node] at (2, 2.2) {};

            % ================= 第二周期: l_N 到 2l_N =================
            % 起点是 m1 (2, 0.8)
            % 假设在 2l_N 时刻，粒子分布在 [1.8, 3.2] 之间
            \fill[cloud-tri] (2, 0.8) -- (4, 1.8) -- (4, 3.2) -- cycle;

            % 高频分支
            \draw[edge-live] (2, 0.8) -- (2.5, 1.1) -- (3.0, 1.4) -- (3.5, 1.3) -- (4, 1.8);
            \draw[edge-killed] (2.5, 1.1) -- (3.2, 2.0) -- (4, 2.3);
            \draw[edge-killed] (3.0, 1.4) -- (3.6, 2.1) -- (4, 2.8);
            \draw[edge-killed] (2, 0.8) -- (2.8, 2.2) -- (4, 3.2);

            % 瓶颈处粒子
            \node[survivor] (m2) at (4, 1.8) {};
            \node[killed-node] at (4, 2.3) {};
            \node[killed-node] at (4, 2.8) {};
            \node[killed-node] at (4, 3.2) {};

            % ================= 第三周期: 2l_N 到 3l_N =================
            % 起点是 m2 (4, 1.8)
            % 假设在 3l_N 时刻，粒子分布在 [2.7, 4.1] 之间
            \fill[cloud-tri] (4, 1.8) -- (6, 2.7) -- (6, 4.1) -- cycle;

            % 高频分支
            \draw[edge-live] (4, 1.8) -- (4.5, 2.0) -- (5.0, 2.4) -- (5.5, 2.3) -- (6, 2.7);
            \draw[edge-killed] (4.5, 2.0) -- (5.2, 2.9) -- (6, 3.2);
            \draw[edge-killed] (5.0, 2.4) -- (5.6, 3.1) -- (6, 3.7);
            \draw[edge-killed] (4, 1.8) -- (4.8, 3.2) -- (6, 4.1);

            % 瓶颈处粒子
            \node[survivor] (m3) at (6, 2.7) {};
            \node[killed-node] at (6, 3.2) {};
            \node[killed-node] at (6, 3.7) {};
            \node[killed-node] at (6, 4.1) {};

            % 瓶颈标注
            \node[above right] at (m1) {\scriptsize keep leftmost};
            \node[above right] at (m2) {\scriptsize keep leftmost};
            \node[above right] at (m3) {\scriptsize keep leftmost};
\end{tikzpicture}

}

\vspace{4pt}
\begin{itemize}
    \item After $\tau_1$, suppress $u_1$, evolve from $v_1$ with $\Xi^{(M)}$
    up to $\tau_1+\ell$, $\ell:=\lfloor(\alpha/\log m)\log N\rfloor$.
    \item \textbf{Iterate the restart functional} $\Phi_{\mathrm{re}}$:
    $\kappa$ = \#restarts until success ($v_\kappa$'s descendants at $\tau_\kappa+\ell$
    are $\ge N^\alpha$ and $\le N$), $\tau:=\tau_\kappa+\ell$ = total time;
    else $\kappa:=\infty$, $\tau:=\infty$.
    \item W.h.p., $\kappa\le K:=\lfloor\delta\log N\rfloor$,
    $\tau\le T:=\gamma\log N$ for some $\delta,\gamma>0$.
\end{itemize}

\end{frame}

\begin{frame}{Upper bound: decomposition and structure}
    % Purpose: define the good event and derive the three-term decomposition.
    % Unified upper bound slide 1: Good event + three-term decomposition (no labels).
% Replaces upper_restart_decomp.tex + upper_term1.tex + upper_term2.tex + upper_term3.tex.

\small
\begin{block}{The ``good'' event}
    $E_N:=\{\kappa\le K,\;\tau\le T\}$, \quad
    $K:=\lfloor\delta\log N\rfloor$, \quad
    $T:=\gamma\log N$.
\end{block}

\vspace{-1mm}
The restart coupling ensures $\mathcal{M}_{\ell_N}^{N|0}\le\mathcal{M}_{\ell_N}^{\mathrm{b}}$ (Lemma~4.8):
\[
    \mathbb{E}\!\left[\mathcal{M}_{\ell_N}^{N|0}\right]
    \le
    \mathbb{E}\!\left[\mathcal{M}_{\ell_N}^{\mathrm{b}}\right]
    =
    \underbrace{\mathbb{E}\!\left[
        \bigl(\mathcal{M}_{\ell_N}^{\mathrm{b}}-\max\mathcal{V}^{\mathrm{b}}_\tau\bigr)
        \mathbf{1}_{E_N}
    \right]}_{\text{Post-restart evolution}}
    +
    \underbrace{\mathbb{E}\!\left[
        \max\mathcal{V}^{\mathrm{b}}_\tau\mathbf{1}_{E_N}
    \right]}_{\text{Restart accumulation}}
    +
    \underbrace{\mathbb{E}\!\left[
        \mathcal{M}_{\ell_N}^{\mathrm{b}}\mathbf{1}_{E_N^{\mathrm{c}}}
    \right]}_{\text{Bad event}}.
\]

\vspace{2mm}
\textbf{Intuition:}
\begin{itemize}
    \item \textbf{Post-restart evolution}: on $E_N$, after $\tau$ the process is a translated $\mathcal{V}^{N|\alpha}$;
    \item \textbf{Restart accumulation}: the maximal position reached during restarts;
    \item \textbf{Bad event}: the contribution when $\kappa>K$ or $\tau>T$, which we will show is negligible.
\end{itemize}

\end{frame}

\begin{frame}{Upper bound: post-restart \& restart accumulation}
    % Purpose: bound the first two terms via Proposition 3.1 and Wald's identity.
    % Unified upper bound slide 2: Bound the post-restart and restart accumulation terms.
% Replaces upper_term1.tex + upper_term2.tex (no equation labels).

\small
\textbf{1. Post-restart evolution.}
On $E_N$, $\tau\le T$ and the process after $\tau$ is a translated $\mathcal{V}^{N|\alpha}$:
\[
    \mathbb{E}\!\left[
        \bigl(\mathcal{M}_{\ell_N}^{\mathrm{b}}-\max\mathcal{V}^{\mathrm{b}}_\tau\bigr)\mathbf{1}_{E_N}
    \right]
    =
    \mathbb{E}\!\left[
        \bigl.\mathbb{E}\!\left[\mathcal{M}^{N|\alpha}_{\ell_N-\iota}\right]\bigr|_{\iota=\tau}\mathbf{1}_{E_N}
    \right]
    \le
    \left(-\alpha+\frac{\pi^2}{2}\sigma^2 t\right)\log N+o_N(\log N),
\]
where the last bound uses \textbf{Theorem~1.1} (the $\mathrm{L}^2$ limit theorem) for $a=\alpha\in(0,1)$.

\vspace{3mm}
\textbf{2. Restart accumulation.}
Define the step sums during restarts:
\[
    G_j:=\tau_j-\tau_{j-1}-1
    \;\sim\;
    \mathrm{Geom}(p_0),
    \qquad
    X_k:=\sum_{i=\tau_{k-1}+1}^{\tau_k-1}
    \bigl(\mathcal{V}^{\mathrm b}(w_{i+1})-\mathcal{V}^{\mathrm b}(w_i)\bigr),
    \qquad
    S_k:=\sum_{i=1}^k X_i,
\]
where $G_j$ (i.i.d.) is the \textbf{time spent} in the $j$-th reboot (until a two-child event occurs),
and $X_k$ is the \textbf{total displacement} during the $k$-th reboot
(a sum of $G_k$ i.i.d.\ copies of $\Xi_1$, the leftmost step).

From the construction, $\max\mathcal{V}^{\mathrm{b}}_\tau\le S_\kappa+(\kappa-1)M+\ell M$, and on $E_N$,
\[
    \mathbb{E}\!\left[\max\mathcal{V}^{\mathrm{b}}_\tau\mathbf{1}_{E_N}\right]
    \le
    \mathbb{E}[S_\kappa]+O_N(\log N).
\]

Since $\kappa$ is a stopping time w.r.t.\ i.i.d.\ step sums, \textbf{Wald's identity} gives
\[
    \mathbb{E}[S_\kappa]
    =
    \mathbb{E}[\kappa]\,\mathbb{E}[X_1]
    =
    \mathbb{E}[\kappa]\,\mathbb{E}[G_1]\,\mathbb{E}[\Xi_1]
    \le
    \frac{1}{p_1}\cdot\frac{1-p_0}{p_0}\cdot\bigl(\mathbb{E}[(\Xi_1)_+]+\mathbb{E}[(\Xi_1)_-]\bigr)
    =O_N(1).
\]

Thus the restart accumulation term is \textbf{bounded}, contributing only $O_N(\log N)$.

\end{frame}

\begin{frame}{Upper bound: the bad event $E_N^{\mathrm{c}}$}
    % Purpose: explain the two-part proof of P(E_N^c) ≤ N^{-ε}.
    % Unified upper bound slide 3: Why P(E_N^c) ≤ N^{-ε}.

\small
\textbf{3. Bad event bound: $\mathbb{P}(E_N^{\mathrm{c}})\le N^{-\varepsilon}$.}
\[
    \mathbb{P}(E_N^{\mathrm{c}})
    \le
    \underbrace{\mathbb{P}(\kappa>K)}_{\text{Too many restarts}}
    \;+\;
    \underbrace{\mathbb{P}(\kappa\le K,\;\tau>T)}_{\text{Restarts take too long}}.
\]

\textbf{Part 1.}
The GW process $\#\mathcal Z^{(M)}$ with offspring law $\Xi^{(M)}$ ($m>1$) has martingale
$\#\mathcal Z^{(M)}_n/m^n\to W$. Since $\#\Xi^{(M)}\le 2$, the $L\log L$ condition holds
automatically, so \textbf{Kesten--Stigum} gives $\mathbb{E}[W]=1$ and $p_1:=\mathbb{P}(W\ge 1)>0$.
Each restart succeeds with probability $\ge p_1$ independently, hence
$\mathbb{P}(\kappa>K)\le(1-p_1)^{\delta\log N}\le N^{-p_1\delta}$.

\textbf{Part 2.}
$G_j:=\tau_j-\tau_{j-1}-1\sim\mathrm{Geom}(p_0)$ i.i.d.\ with success probability
$p_0=\mathbb{P}(\#\Xi^{(M)\mathrm b}\ge 2)>0$.
Since $G_j$ is geometric, it has a finite exponential moment:
for any $0<\theta<-\log(1-p_0)$,
$\mathbb{E}[e^{\theta G_1}]=\frac{p_0}{1-(1-p_0)e^\theta}<\infty$.
Then $\tau=\tau_\kappa+\ell$ with $\tau_k=\sum_{j=1}^k(G_j+1)$,
and by \textbf{exponential Markov} (Chernoff bound),
\[
    \mathbb{P}(\kappa\le K,\tau>T)
    \le\frac{\mathbb{E}[e^{\theta G_1}]^K}{e^{\theta(T-K-\ell)}}
    =N^{\theta(\delta+\frac{\alpha}{\log m}-\gamma)+\delta\log\mathbb{E}[e^{\theta G_1}]}
    =:N^{-\varepsilon_2},
\]
where $\gamma>0$ is chosen large enough so that $\varepsilon_2>0$.

\textbf{Contribution to the upper bound.}
Cauchy--Schwarz gives
$|\mathbb{E}[\mathcal{M}_{\ell_N}^{\mathrm{b}}\mathbf{1}_{E_N^{\mathrm{c}}}]|
\le\mathbb{E}[(\mathcal{M}_{\ell_N}^{\mathrm{b}})^2]^{1/2}\mathbb{P}(E_N^{\mathrm{c}})^{1/2}
\le O(\log^3 N)\cdot N^{-\varepsilon/2}\to 0$.

\textbf{Note:} the $\mathbb{E}[(\Xi_1)_+^4]<\infty$ assumption used here for the $O(\log^3 N)$
moment bound is \textbf{no longer needed} after a later refinement;
only $\mathbb{E}[(\Xi_1)_+]<\infty$ suffices.

\end{frame}

\begin{frame}{Upper bound: conclusion}
    % Purpose: assemble the three terms and deduce the final bound via t→∞.
    \small
    From the three parts ($\alpha$ arbitrarily small):
    \[
        \mathbb{E}\!\left[\mathcal{M}_{\ell_N}^{\mathrm{b}}\right]
        \le
        \frac{\pi^2}{2}\sigma^2 t\log N + O_N(\log N).
    \]

    By the coupling $\mathcal{M}_{\ell_N}^{N|0}\le\mathcal{M}_{\ell_N}^{\mathrm{b}}$ and
    $\ell_N = t\log^3 N + o_N(\log^3 N)$,
    \begin{align*}
        v_N
        &\le \frac{1}{\ell_N}\,\mathbb{E}\!\left[\mathcal{M}_{\ell_N}^{\mathrm{b}}\right] \\[3pt]
        &\le \frac{1}{t\log^3 N + o_N(\log^3 N)}
            \Bigl[\,\frac{\pi^2}{2}\sigma^2 t\log N + O_N(\log N)\Bigr] \\[3pt]
        &\le \inf_{t>0}\Bigl[
              \Bigl(\frac{O_N(1)}{t} + \frac{\pi^2}{2}\sigma^2\Bigr)\frac{1}{\log^2 N}
              + o_N\Bigl(\frac{1}{\log^2 N}\Bigr)
            \Bigr] \\[3pt]
        &= \frac{\pi^2}{2}\sigma^2\frac{1}{\log^2 N}
           + o_N\!\left(\frac{1}{\log^2 N}\right).
    \end{align*}
    The $\inf_{t>0}$ sends the first term to $0$, leaving the $\frac{\pi^2}{2}\sigma^2$ constant.
\end{frame}

% ==========================================
\section{Lower bound: proof from above}
% ==========================================

\begin{frame}{Lower comparison construction}
    % Purpose: construction and illustration (from contents).
    \begin{columns}
        \column{0.35\textwidth}
        \small
        % Purpose: define the lower comparison process used for the speed lower bound.
\small
At time $0$, raise $\lfloor N^\alpha\rfloor$ to $N^1$ particles
(keep leftmost $N$ among all descendants); the original process
is embedded as a subtree of this new process.
Evolve as $\mathcal V^{N|1}$; at each $k\ell_N$,
translate all particles to the leftmost position.

        \column{0.65\textwidth}
        \centering
        \resizebox{\textwidth}{!}{% Source: contents/part-5-eng.tex (Figure: construction of V^{lower})
\begin{tikzpicture}[
    xscale=2,
    yscale=1.2,
    >=Stealth,
    % styles
    survivor/.style={circle, fill=blue!70!black, inner sep=1.2pt},
    other-node/.style={circle, draw=blue!70!black, fill=white, inner sep=1pt},
    edge-live/.style={thick, blue!70!black},
    edge-killed/.style={thick, dashed, blue!30},
    cloud-tri/.style={fill=blue!10, opacity=0.5},
    translation-arrow/.style={->, thick, red!80!black, dashed}
]

    % reference times (0, l_N, 2l_N, 3l_N)
    \foreach \x/\label in {0/0, 2/\ell_N, 4/2\ell_N, 6/3\ell_N} {
        \draw[gray!40, dashed] (\x, -0.5) -- (\x, 4.5);
        \node[below] at (\x, -0.5) {$\label$};
    }

    % ===== Block 1: 0 -> l_N
    \node[survivor] (n0) at (0, 0) {};
    \fill[cloud-tri] (0, 0) -- (2, 0.8) -- (2, 2.2) -- cycle;
    \draw[edge-live] (0,0) -- (0.8, 0.3) -- (1.4, 0.6) -- (2, 0.8);
    \draw[edge-live] (0.8, 0.3) -- (1.3, 1.2) -- (2, 1.3);
    \draw[edge-live] (1.4, 0.6) -- (1.8, 1.5) -- (2, 1.8);
    \draw[edge-live] (0,0) -- (0.9, 1.2) -- (2, 2.2);
    \node[survivor] (m1) at (2, 0.8) {};
    \node[survivor] at (2, 1.3) {};
    \node[survivor] at (2, 1.8) {};
    \node[survivor] at (2, 2.2) {};

    % translations at time l_N
    \draw[translation-arrow] (2, 2.2) to[bend left] (2, 0.9);
    \draw[translation-arrow] (2, 1.8) to[bend left] (2, 0.85);
    \draw[translation-arrow] (2, 1.3) to[bend left] (2, 0.82);

    % ===== Block 2: l_N -> 2l_N (all N particles start from m1)
    \fill[cloud-tri] (2, 0.8) -- (4, 1.8) -- (4, 3.2) -- cycle;
    \draw[edge-live] (2, 0.8) -- (2.5, 1.1) -- (3.0, 1.6) -- (3.5, 1.5) -- (4, 1.8);
    \draw[edge-live] (2, 0.8) -- (2.6, 1.8) -- (3.2, 2.2) -- (4, 2.5);
    \draw[edge-live] (2, 0.8) -- (2.7, 2.1) -- (3.6, 2.9) -- (4, 3.2);
    \node[survivor] (m2) at (4, 1.8) {};
    \node[survivor] at (4, 2.5) {};
    \node[survivor] at (4, 3.2) {};

    \draw[translation-arrow] (4, 3.2) to[bend left] (4, 1.9);
    \draw[translation-arrow] (4, 2.5) to[bend left] (4, 1.85);

    % ===== Block 3: 2l_N -> 3l_N
    \fill[cloud-tri] (4, 1.8) -- (6, 2.7) -- (6, 4.1) -- cycle;
    \draw[edge-live] (4, 1.8) -- (4.5, 2.1) -- (5.1, 2.5) -- (6, 2.7);
    \draw[edge-live] (4, 1.8) -- (4.8, 2.9) -- (5.5, 3.4) -- (6, 4.1);
    \draw[edge-live] (4, 1.8) -- (4.6, 2.5) -- (5.3, 3.0) -- (6, 3.4);
    \node[survivor] (m3) at (6, 2.7) {};
    \node[survivor] at (6, 3.4) {};
    \node[survivor] at (6, 4.1) {};

    % labels
    \node[below, blue!70!black] at (2, 0.8) {$\mathcal{M}^N_{\ell_N}$};
    \node[below, blue!70!black] at (4, 1.8) {$\mathcal{M}^N_{2\ell_N}$};
    \node[below, blue!70!black] at (6, 2.7) {$\mathcal{M}^N_{3\ell_N}$};
\end{tikzpicture}
}
    \end{columns}
\end{frame}

\begin{frame}{Lower bound: SLLN + Theorem~1.1}
    % Purpose: simplified chain with SLLN and Theorem 1.1 explicit.
    % Simplified lower bound inequality chain with SLLN and Theorem 1.1.
\small
\begin{align*}
    v_N
    =
    \lim_{k\to\infty}
    \frac 1{k\ell_N}
    \mathcal M^{N|a}_{k\ell_N}
    &\ge
    \frac 1{\ell_N}
    \lim_{k\to\infty}
    \frac 1k
    \sum_{i=1}^k
    \min\bigl(
        \mathcal V^{\mathrm{lower}}_{i\ell_N}
        -
        \mathcal M^{\mathrm{lower}}_{(i-1)\ell_N}
    \bigr)
    \\
    &\stackrel{\text{SLLN}}{=}
    \frac 1{\ell_N}\,
    \mathbb E\!\left[\mathcal M^{N|1}_{\ell_N}\right]
    \\
    &\stackrel{\text{Thm~1.1}}{=}
    \frac 1{t\log^3N+o(\log^3N)}
    \Bigl(
        \bigl(-1+\tfrac{\pi^2}2\sigma^2t\bigr)\log N
        +
        o_N(\log N)
    \Bigr)
    \\
    &=
    \sup_{t>0}\,
    \Bigl(-\frac 1t+\frac{\pi^2}2\sigma^2\Bigr)\frac{1}{\log^2 N}
    + o_N\!\left(\frac{1}{\log^2 N}\right)
    \\[3pt]
    &=
    \frac{\pi^2}{2}\sigma^2\frac{1}{\log^2 N}
    + o_N\!\left(\frac{1}{\log^2 N}\right).
\end{align*}
The $\sup_{t>0}$ sends $-1/t\to 0$.  SLLN applies because the increments
$\min(\mathcal V^{\mathrm{lower}}_{i\ell_N}-\mathcal M^{\mathrm{lower}}_{(i-1)\ell_N})$
are i.i.d.\ with finite mean $\mathbb E[\mathcal M^{N|1}_{\ell_N}]<\infty$ (by
$\mathbb E[(\Xi_1)_+]<\infty$).  The limit constant is given by
\textbf{Theorem~1.1} ($\mathrm L^2$ version, $a=1$; can be refined to
$\mathrm L^1$).

\end{frame}

% (Moved/removed) Standard inputs and tube picture are introduced later where needed (Theorem 1.1 proof).

% ==========================================
\section{The $\mathrm L^2$ limit theorem and its lemma chain}
% ==========================================

\begin{frame}{Theorem 1.1 + $\mathrm L^2$ reduction (one slide)}
    % Purpose: restate Theorem 1.1 (incl. a=0) and reduce L^2 to two tail bounds, concisely.
    \small
    % Theorem 1.1 (restated, including the a=0 extension condition).
\begin{theorem}[Theorem 1.1]
Fix $a\in(0,1]$ and $t>0$. If $\mathcal V^N=\mathcal V^{N|a}$ starts with $\lfloor N^a\rfloor$ particles located at $0$, and its branching law $\Xi$ satisfies the standing assumptions in the thesis (boundary case, finite variance, finite cross second moment, and moment conditions on $\Xi_1$), then for any sequence $\ell_N$ satisfying $\lim_{N\to\infty}\ell_N/\log^3N=t$,
\[
    \frac 1{\log N}\min_{|u|=\ell_N}\mathcal V^N(u)
    \xrightarrow{\mathrm L^2}
    -a+\frac{\pi^2}{2}\sigma^2 t.
\]
\vspace{1mm}
 \textbf{Moreover (the $a=0$ version).} Under the additional exponential-moment condition on the leftmost step (Assumption 9 in the thesis), the same $\mathrm L^2$ convergence also holds for $a=0$.
\end{theorem}

    \vspace{-1mm}
    Let $\mathcal X_N:=\mathcal M^N_{\ell_N}/\log N$ and $\mu:=-a+\frac{\pi^2}{2}\sigma^2 t$.
    Fix $\zeta<\mu<\zeta'$ and set $H_N:=\{\zeta\le \mathcal X_N\le \zeta'\}$.
    Using the 4th-moment control of $\mathcal M^N_n$ and Cauchy--Schwarz,
    \[
        \E[\mathcal X_N^2\mathbf 1_{H_N^{\mathrm c}}]
        \le
        \E[\mathcal X_N^4]^{1/2}\,\P(H_N^{\mathrm c})^{1/2},
    \]
    hence $\mathcal X_N\xrightarrow{\mathrm L^2}\mu$ follows once we prove:
    \[
        \P(\mathcal M^N_{\ell_N}\ge \zeta'\log N)\le N^{-\varepsilon},
        \qquad
        \P(\mathcal M^N_{\ell_N}\le \zeta\log N)\le N^{-\varepsilon}.
    \]
\end{frame}

\begin{frame}{Right tail upper bound, $a\in(0,1)$: killed BRW construction (A\"id\'ekon--Hu)}
    % Purpose: introduce the killed BRW and state the phase-by-phase killing rule.
    \small
    \begin{center}
        \resizebox{!}{0.56\textheight}{\begin{tikzpicture}[
        xscale=1.75, 
        yscale=0.9, 
        >=Stealth,
        % 样式定义
        survivor/.style={circle, fill=blue!70!black, inner sep=1.2pt},
        other-node/.style={circle, draw=blue!70!black, fill=white, inner sep=1pt},
        killed-node/.style={circle, draw=blue!70!black, fill=white, inner sep=1pt, cross out, thick, minimum size=3pt},
        edge-live/.style={thick, blue!70!black},
        edge-killed/.style={thick, dashed, blue!70!black}, 
        tube-fill/.style={blue!10, opacity=0.4},
        tube-edge/.style={thin, black},
        ref-dash/.style={thin, gray!60, dashed}
    ]

        % 1. 坐标轴
        \draw[->, thick] (0, -2.5) -- (0, 3.8) node[above] {$\mathcal{V}(u)$};
        \draw[->, thick, dashed] (0, 0) -- (6.5, 0) node[right] {$n$};

        % 2. 阶段 I: 水平管道 [0, 1.5]
        \fill[tube-fill] (0, -2.0) rectangle (1.5, 0);
        \draw[tube-edge] (0, -2.0) -- (1.5, -2.0) node[pos=0, left, black] {$-b \log N$};
        \draw[tube-edge] (0, 0) -- (1.5, 0);

        % 3. 阶段 II: 倾斜管道 [1.5, 5.0]
        \def\ystart{-1.8} 
        \def\ylow{1.2}  % 终点下界
        \def\ywin{1.8}  % 终点窗口界
        \def\yhigh{2.8} % 终点上界
        \fill[tube-fill] (1.5, \ystart) -- (5.0, \ylow) -- (5.0, \yhigh) -- (1.5, \ystart+1.6) -- cycle;
        \draw[tube-edge] (1.5, \ystart) -- (5.0, \ylow); 
        \draw[tube-edge] (1.5, \ystart+1.6) -- (5.0, \yhigh); 
        
        \node[survivor] (y-node) at (1.5, \ystart) {};
        \node[left, xshift=-2pt] at (1.5, \ystart) {$y$};

        % 4. 均匀高频分支 (步长 0.3)
        \node[survivor] (n0) at (0, 0) {};

        % --- 阶段 I 路径 ---
        \foreach \i/\x/\y in {1/0.3/-0.4, 2/0.6/-0.1, 3/0.9/-0.8, 4/1.2/-1.4} {
            \coordinate (p\i) at (\x, \y);
            \node[survivor] at (p\i) {};
        }
        \draw[edge-live] (n0) -- (p1) -- (p2) -- (p3) -- (p4) -- (y-node);
        
        % 阶段 I 杀掉的分支 (管道外)
        \draw[edge-killed] (p1) -- (0.6, 0.4); \node[killed-node] at (0.6, 0.4) {};
        \draw[edge-killed] (p3) -- (1.2, -2.3); \node[killed-node] at (1.2, -2.3) {};
        \draw[edge-killed] (p4) -- (1.5, -1); \node[killed-node] at (1.5, -1) {};

        % --- 阶段 II 路径 (向上移动) ---
        \foreach \i/\x/\y in {1/1.8/-1.5, 2/2.1/-1.2, 3/2.4/-0.8, 4/2.7/-0.5, 5/3.0/-0.1, 6/3.3/0.3, 7/3.6/0.6, 8/3.9/0.9, 9/4.2/1.1, 10/4.5/1.3, 11/4.8/1.5} {
            \coordinate (q\i) at (\x, \y);
            \node[survivor] at (q\i) {};
        }
        \coordinate (m-node) at (5.0, 1.6);
        \draw[edge-live] (y-node) -- (q1) -- (q2) -- (q3) -- (q4) -- (q5) -- (q6) -- (q7) -- (q8) -- (q9) -- (q10) -- (q11) -- (m-node);
        \node[survivor] at (m-node) {};

        % 阶段 II 存活但非极小的分支 (管道内)
        \draw[edge-live] (q6) -- (3.6, 1.1) -- (3.9, 1.5) -- (4.2, 1.8) -- (4.5, 2.1) -- (4.8, 2.3);
        \foreach \x/\y in {3.6/1.1, 3.9/1.5, 4.2/1.8, 4.5/2.1, 4.8/2.3} \node[other-node] at (\x, \y) {};

        % 阶段 II 杀掉的分支 (越界)
        \draw[edge-killed] (q9) -- (4.5, 0.3); \node[killed-node] at (4.5, 0.3) {};
        \draw[edge-killed] (4.2, 1.8) -- (4.5, 2.8); \node[killed-node] at (4.5, 2.8) {};
        \draw[edge-killed] (4.8, 2.3) -- (5.0, 2.1); \node[killed-node] at (5.0, 2.1) {};

        % 5. 时间节点辅助线
        \draw[dashed, black!30] (1.5, -2.5) -- (1.5, 3.5) node[above, black] {$\ell$};
        \draw[dashed, black!30] (5.0, -2.5) -- (5.0, 3.5) node[above, black] {$\ell_N$};

        % 6. 标注窗口与跨度 (向右侧引出)
        
        % 左侧时刻 l 窗口
        \draw[ultra thick, blue!80!black] (1.5, -2.0) -- (1.5, -1.3);
        \draw[decorate, decoration={brace, amplitude=4pt, mirror}] (2, -2.0) -- (2, -1.3) 
            node[midway, right, xshift=2pt] {$v \log N$};

        % --- 右侧引出线与标注 ---
        % 引出辅助水平线
        \draw[ref-dash] (5.0, \ylow) -- (6.2, \ylow);
        \draw[ref-dash] (5.0, \ywin) -- (5.8, \ywin);
        \draw[ref-dash] (5.0, \yhigh) -- (6.2, \yhigh);
        \draw[ref-dash] (1.53, -2.0) -- (2, -2.0);
        \draw[ref-dash] (1.53, -1.3) -- (2, -1.3);

        % 标注 M^k (紧贴 L_N 垂线右侧)
        \node[right, blue!70!black] at (5, 1.6) {$\mathcal{M}^{\mathrm{k}}_{\ell_N}$};

        % 标注 vL^1/3 (第一层大括号)
        \draw[ultra thick, blue!80!black] (5, \ylow) -- (5, \ywin);
        \draw[decorate, decoration={brace, amplitude=4pt, mirror}] (5.6, \ylow) -- (5.6, \ywin) 
            node[midway, right, xshift=1pt] {$vL^{1/3}$};

        % 标注 cL^1/3 (第二层更远的大括号)
        \draw[decorate, decoration={brace, amplitude=5pt, mirror}] (6.3, \ylow) -- (6.3, \yhigh) 
            node[midway, right, xshift=2pt] {$cL^{1/3}$};

        % 7. 阈值线
        \draw[red!80!black, thick, dashed] (0, 3.2) -- (5.5, 3.2) node[right] {$\zeta' \log N$};

    \end{tikzpicture}
}
    \end{center}
    \vspace{-2mm}
    \footnotesize
    Replace the selection-based $N$-BRW by a \emph{killed} BRW $\mathcal V^{\mathrm k}$ (tube constraints).
    Let $\ell:=\lfloor\log^{5/2}N\rfloor$ and $L:=\ell_N-\ell$.
    \vspace{-1mm}
    \begin{itemize}
        \item Phase I ($0\to \ell$): kill outside $[-b\log N,0]$, then keep only $[-b\log N,(-b+v)\log N]$ at time $\ell$.
        \item Phase II ($\ell+1\to \ell_N$): for each time-$\ell$ ancestor at $y$, kill outside $y+[r_j,s_j]$ with
        $r_j:=\frac{\pi^2}{2}\sigma^2c^{-2}L^{-2/3}(j-\ell)$ and $s_j:=r_j+cL^{1/3}$, then keep only $y+[r_{\ell_N},\,r_{\ell_N}+vL^{1/3}]$.
    \end{itemize}
\end{frame}

\begin{frame}{Killed BRW construction: parameters and ranges (A\"id\'ekon--Hu)}
    % Purpose: record admissible parameter ranges; keep notation identical to the thesis.
    \small
    \begin{center}
        \resizebox{!}{0.52\textheight}{\begin{tikzpicture}[
        xscale=1.75, 
        yscale=0.9, 
        >=Stealth,
        % 样式定义
        survivor/.style={circle, fill=blue!70!black, inner sep=1.2pt},
        other-node/.style={circle, draw=blue!70!black, fill=white, inner sep=1pt},
        killed-node/.style={circle, draw=blue!70!black, fill=white, inner sep=1pt, cross out, thick, minimum size=3pt},
        edge-live/.style={thick, blue!70!black},
        edge-killed/.style={thick, dashed, blue!70!black}, 
        tube-fill/.style={blue!10, opacity=0.4},
        tube-edge/.style={thin, black},
        ref-dash/.style={thin, gray!60, dashed}
    ]

        % 1. 坐标轴
        \draw[->, thick] (0, -2.5) -- (0, 3.8) node[above] {$\mathcal{V}(u)$};
        \draw[->, thick, dashed] (0, 0) -- (6.5, 0) node[right] {$n$};

        % 2. 阶段 I: 水平管道 [0, 1.5]
        \fill[tube-fill] (0, -2.0) rectangle (1.5, 0);
        \draw[tube-edge] (0, -2.0) -- (1.5, -2.0) node[pos=0, left, black] {$-b \log N$};
        \draw[tube-edge] (0, 0) -- (1.5, 0);

        % 3. 阶段 II: 倾斜管道 [1.5, 5.0]
        \def\ystart{-1.8} 
        \def\ylow{1.2}  % 终点下界
        \def\ywin{1.8}  % 终点窗口界
        \def\yhigh{2.8} % 终点上界
        \fill[tube-fill] (1.5, \ystart) -- (5.0, \ylow) -- (5.0, \yhigh) -- (1.5, \ystart+1.6) -- cycle;
        \draw[tube-edge] (1.5, \ystart) -- (5.0, \ylow); 
        \draw[tube-edge] (1.5, \ystart+1.6) -- (5.0, \yhigh); 
        
        \node[survivor] (y-node) at (1.5, \ystart) {};
        \node[left, xshift=-2pt] at (1.5, \ystart) {$y$};

        % 4. 均匀高频分支 (步长 0.3)
        \node[survivor] (n0) at (0, 0) {};

        % --- 阶段 I 路径 ---
        \foreach \i/\x/\y in {1/0.3/-0.4, 2/0.6/-0.1, 3/0.9/-0.8, 4/1.2/-1.4} {
            \coordinate (p\i) at (\x, \y);
            \node[survivor] at (p\i) {};
        }
        \draw[edge-live] (n0) -- (p1) -- (p2) -- (p3) -- (p4) -- (y-node);
        
        % 阶段 I 杀掉的分支 (管道外)
        \draw[edge-killed] (p1) -- (0.6, 0.4); \node[killed-node] at (0.6, 0.4) {};
        \draw[edge-killed] (p3) -- (1.2, -2.3); \node[killed-node] at (1.2, -2.3) {};
        \draw[edge-killed] (p4) -- (1.5, -1); \node[killed-node] at (1.5, -1) {};

        % --- 阶段 II 路径 (向上移动) ---
        \foreach \i/\x/\y in {1/1.8/-1.5, 2/2.1/-1.2, 3/2.4/-0.8, 4/2.7/-0.5, 5/3.0/-0.1, 6/3.3/0.3, 7/3.6/0.6, 8/3.9/0.9, 9/4.2/1.1, 10/4.5/1.3, 11/4.8/1.5} {
            \coordinate (q\i) at (\x, \y);
            \node[survivor] at (q\i) {};
        }
        \coordinate (m-node) at (5.0, 1.6);
        \draw[edge-live] (y-node) -- (q1) -- (q2) -- (q3) -- (q4) -- (q5) -- (q6) -- (q7) -- (q8) -- (q9) -- (q10) -- (q11) -- (m-node);
        \node[survivor] at (m-node) {};

        % 阶段 II 存活但非极小的分支 (管道内)
        \draw[edge-live] (q6) -- (3.6, 1.1) -- (3.9, 1.5) -- (4.2, 1.8) -- (4.5, 2.1) -- (4.8, 2.3);
        \foreach \x/\y in {3.6/1.1, 3.9/1.5, 4.2/1.8, 4.5/2.1, 4.8/2.3} \node[other-node] at (\x, \y) {};

        % 阶段 II 杀掉的分支 (越界)
        \draw[edge-killed] (q9) -- (4.5, 0.3); \node[killed-node] at (4.5, 0.3) {};
        \draw[edge-killed] (4.2, 1.8) -- (4.5, 2.8); \node[killed-node] at (4.5, 2.8) {};
        \draw[edge-killed] (4.8, 2.3) -- (5.0, 2.1); \node[killed-node] at (5.0, 2.1) {};

        % 5. 时间节点辅助线
        \draw[dashed, black!30] (1.5, -2.5) -- (1.5, 3.5) node[above, black] {$\ell$};
        \draw[dashed, black!30] (5.0, -2.5) -- (5.0, 3.5) node[above, black] {$\ell_N$};

        % 6. 标注窗口与跨度 (向右侧引出)
        
        % 左侧时刻 l 窗口
        \draw[ultra thick, blue!80!black] (1.5, -2.0) -- (1.5, -1.3);
        \draw[decorate, decoration={brace, amplitude=4pt, mirror}] (2, -2.0) -- (2, -1.3) 
            node[midway, right, xshift=2pt] {$v \log N$};

        % --- 右侧引出线与标注 ---
        % 引出辅助水平线
        \draw[ref-dash] (5.0, \ylow) -- (6.2, \ylow);
        \draw[ref-dash] (5.0, \ywin) -- (5.8, \ywin);
        \draw[ref-dash] (5.0, \yhigh) -- (6.2, \yhigh);
        \draw[ref-dash] (1.53, -2.0) -- (2, -2.0);
        \draw[ref-dash] (1.53, -1.3) -- (2, -1.3);

        % 标注 M^k (紧贴 L_N 垂线右侧)
        \node[right, blue!70!black] at (5, 1.6) {$\mathcal{M}^{\mathrm{k}}_{\ell_N}$};

        % 标注 vL^1/3 (第一层大括号)
        \draw[ultra thick, blue!80!black] (5, \ylow) -- (5, \ywin);
        \draw[decorate, decoration={brace, amplitude=4pt, mirror}] (5.6, \ylow) -- (5.6, \ywin) 
            node[midway, right, xshift=1pt] {$vL^{1/3}$};

        % 标注 cL^1/3 (第二层更远的大括号)
        \draw[decorate, decoration={brace, amplitude=5pt, mirror}] (6.3, \ylow) -- (6.3, \yhigh) 
            node[midway, right, xshift=2pt] {$cL^{1/3}$};

        % 7. 阈值线
        \draw[red!80!black, thick, dashed] (0, 3.2) -- (5.5, 3.2) node[right] {$\zeta' \log N$};

    \end{tikzpicture}
}
    \end{center}
    \vspace{-1mm}
    \footnotesize
    Fix $c\in(0,t^{-1/3})$, $b\in(0,a)$ and $v>0$ such that
    \[
        ct^{1/3}+(a-b)+v<1,
        \qquad
        -b+\frac{\pi^2}{2}\sigma^2\!\left(c^{-2}t^{1/3}\right)+2v(1+t^{1/3})<\zeta'.
    \]
    \vspace{1mm}
    Parameters are fixed with the intended regime: $c\downarrow 0$ (narrow tube), $v\downarrow 0$ (thin terminal window), $b\uparrow a$ (initial cutoff close to $-a\log N$).
\end{frame}

\begin{frame}{Coupling + right-tail decomposition}
    \small
    \begin{itemize}
        \item (Lemma 3.17) If the killing rule is adapted and $\mathcal V^{\mathrm k}_0\subseteq \mathcal V^N_0$, then one can couple $(\mathcal V^N,\mathcal V^{\mathrm k})$ so that
        \[
            \mathcal M^N_k\le \mathcal M^{\mathrm k}_k,\qquad
            k\le\tau:=\inf\{k:\#\mathcal Z^{\mathrm k}_k>N\}.
        \]
        \item Hence (right tail, $a\in(0,1)$),
        \[
            \P\!\left(\mathcal M^N_{\ell_N}\ge \zeta'\log N\right)
            \le
            \P(\tau\le \ell_N)
            +
            \P\!\left(\mathcal M^{\mathrm k}_{\ell_N}\ge \zeta'\log N\right).
        \]
    \end{itemize}
\end{frame}

\begin{frame}{Right tail upper bound ($a\in(0,1)$): Step 1 (control $\tau$)}
    \small
    % Right tail Step 1: control P(tau <= ell_N) by first moments + stage-wise population estimates.
\begin{block}{Step 1: bound $\P(\tau\le \ell_N)$ (union bound + Markov)}
\[
    \P(\tau\le \ell_N)
    \le
    \sum_{n=1}^{\ell_N}\P(\#\mathcal Z^{\mathrm k}_n>N)
    \le
    \sum_{n=1}^{\ell_N}\frac{\E[\#\mathcal Z^{\mathrm k}_n]}{N}.
\]
\end{block}
\vspace{-1mm}
\begin{itemize}
    \item Phase I ($1\le n\le \ell$): $\E[\#\mathcal Z^{\mathrm k}_n]\le N^{a+o_N(1)}$ using Eq.~(4.10) in Lemma 3.10.
    \item Bottleneck at time $\ell$:
    $\E[\#\mathcal Z^{\mathrm k}_{\ell+}]=N^{a-b+v+o_N(1)}$ using Eq.~(4.11) in Lemma 3.10.
    \item Phase II ($\ell<n\le \ell_N$): $\E[\#\mathcal Z^{\mathrm k}_n]\le N^{a-b+v+ct^{1/3}+o_N(1)}$ using Eq.~(4.8) in Lemma 3.9 + branching property.
\end{itemize}
\vspace{-1mm}
Thus $\P(\tau\le \ell_N)\le N^{-(1-\eta)+o_N(1)}$ with $\eta:=\max\{a,\ a-b+v+ct^{1/3}\}<1$.

\end{frame}

\begin{frame}{Right tail upper bound ($a\in(0,1)$): Step 2 (control $\mathcal M^{\mathrm k}_{\ell_N}$)}
    \small
    % Right tail Step 2: control P(M^k_{ell_N} >= zeta' log N) via extinction at ell_N+ and a split at ell+.
\begin{block}{Step 2: bound $\P(\mathcal M^{\mathrm k}_{\ell_N}\ge \zeta'\log N)$}
In the construction, any alive particle at time $\ell_N$ stays below $\zeta'\log N$, hence
\[
    \P(\mathcal M^{\mathrm k}_{\ell_N}\ge \zeta'\log N)
    \le
    \P(\#\mathcal Z^{\mathrm k}_{\ell_N+}=0).
\]
\end{block}
\vspace{-1mm}
\begin{block}{Split at time $\ell$}
Fix $\delta\in(0,\tfrac12(a-b+v))$. Then
\[
    \P(\#\mathcal Z^{\mathrm k}_{\ell_N+}=0)
    \le
    \P(\#\mathcal Z^{\mathrm k}_{\ell+}<N^\delta)
    +
    \P(\#\mathcal Z^{\mathrm k}_{\ell_N+}=0,\ \#\mathcal Z^{\mathrm k}_{\ell+}\ge N^\delta).
\]
\end{block}
\vspace{-1mm}
\begin{itemize}
    \item \textbf{Few particles at $\ell+$:} lower-bound $\#\mathcal Z^{\mathrm k}_{\ell+}$ by a binomial variable using Lemma 3.10, then apply Chernoff.
    \item \textbf{Many particles at $\ell+$:} condition on $\mathcal F_\ell$ and use independence of subtrees (one factor per $\nu\in\mathcal Z^{\mathrm k}_{\ell+}$).
\end{itemize}

\end{frame}

\begin{frame}{Right tail upper bound ($a\in(0,1)$): tools used (grouped by parts)}
    % Purpose: list the exact tools used, grouped by where they come from in contents/.
    \small
    \begin{itemize}
        \item \textbf{Coupling:} Lemma 3.17 transfers killed-BRW bounds to the $N$-BRW up to $\tau$.
        \item \textbf{Second moment chain:} Lemma 3.8 (\textbf{uses Assumption 6}) $\Rightarrow$ Lemma 3.9 (\textbf{Paley--Zygmund}) $\Rightarrow$ Lemma 3.10.
        \item \textbf{Spine walk:} Lemma 3.4, Lemma 3.5, Lemma 3.6, Lemma 3.7.
        \item \textbf{Known input:} Theorem 2.4 (Mogul'skii) provides the sharp tube constant used in Lemma 3.4--3.7.
    \end{itemize}
\end{frame}

\begin{frame}{Spine-walk lemma group I: horizontal tube (Lemma 3.4)}
    % Purpose: explains what problem this lemma solves.
    \small
    % Lemma `lemma:4.1` statement (compact, faithful to contents/part-2-eng.tex).
\begin{lemma}[\texttt{lemma:4.1}]
Let $v>0$ and $0<d<\frac{\pi^2}2\sigma^2$. There exists $\Delta_0>0$ such that for all $\Delta\ge\Delta_0$, $r_0\le 0\le r_0+\Delta$ and $L\ge 1$,
\[
    \mathbb P\big(r_0\le Y_j\le r_0+\Delta,\ \forall j\in[\![1,L]\!]\big)
    \le
    \exp\!\left(\Delta^v-\frac{d}{\Delta^2}L\right).
\]
\end{lemma}


    \vspace{1mm}
    \centering
    \resizebox{0.72\textwidth}{!}{\begin{tikzpicture}[
        scale=1.5,
        >=Stealth,
        % 随机游走路径样式
        rw-path/.style={thick, blue!70!black, mark=*, mark size=1.5pt},
        % 管道填充样式
        tube-fill/.style={gray!10},
        % 标注样式
        label-font/.style={font=\small}
    ]

        % 1. 绘制背景管道填充 (r_0 到 r_0 + delta)
        \fill[tube-fill] (0, -1.5) rectangle (4, 0.5);

        % 2. 绘制辅助虚线
        \draw[dashed, black!60] (0, -2) -- (4.5, -2) node[right] {$-\Delta$};
        % \draw[dashed, black!60] (0, 0) -- (4.5, 0) node[right] {$y = 0$};
        \draw[dashed, black!60] (4, -2.5) -- (4, 1) node[above] at (4, 1) {$L$};

        % 3. 绘制管道实线 (r_0 到 r_0 + delta)
        \draw[thick, black] (0, -1.5) -- (4.5, -1.5) node[right] {$r_0$};
        \draw[thick, black] (0, 0.5) -- (4.5, 0.5) node[right] {$r_0 + \Delta$};

        % 4. 绘制坐标轴 (Y轴实线)
        \draw[->, thick] (0, -2.5) -- (0, 1.2) node[above] {$Y_n$};
        \draw[->, thick, dashed] (0, 0) -- (5, 0) node[right] {$n$};

        % 5. 绘制随机游走 (起点在 0,0，始终保持在 [-1.5, 0.5] 之间)
        % 选取了一些看起来随机的点
        \draw[rw-path] plot coordinates {
            (0, 0) 
            (0.8, -0.6) 
            (1.5, 0.2) 
            (2.2, -1.1) 
            (3.1, -0.4) 
            (4, -0.8)
        };

        % 6. 标注 Delta (使用大括号)
        \draw[decorate, decoration={brace, amplitude=5pt}] 
            (-0.1, -1.5) -- (-0.1, 0.5) node[midway, left=6pt] {$\Delta$};

        % 标注特定的点
        % \node[below left] at (0,0) {$0$};

    \end{tikzpicture}}
\end{frame}

\begin{frame}{Spine-walk lemma group II: no large drop (Lemma 3.5)}
    % Purpose: prevents large drawdowns; used to control tube events with extra constraints.
    \small
    % Lemma `lemma:2.3` statement (compact, faithful to contents/part-2-eng.tex).
\begin{lemma}[\texttt{lemma:2.3}]
Let $c>0$. For any $d\in(0,\frac{\pi^2}2\sigma^2)$, there exists $\Delta_0=\Delta_0(c,d)\ge 1$ such that for all $\Delta\ge\Delta_0$ and $L\ge\Delta^{2+c}$,
\[
    \mathbb P\!\left(
        Y_L\le\frac{\Delta}{c},\ 
        \min_{0\le i\le j\le L}(Y_j-Y_i)\ge-\Delta
    \right)
    \le
    e^{-\frac{dL}{\Delta^2}}.
\]
\end{lemma}


    \vspace{1mm}
    \centering
    \resizebox{0.78\textwidth}{!}{\begin{tikzpicture}[
        xscale=1.7, % 修改这里控制横向（代数 n）的间距
        yscale=1.1, % 修改这里控制纵向（空间位置 V）的伸展程度
        >=Stealth,
        % 密集随机游走样式
        rw-path/.style={thick, blue!70!black},
        % 动态屏障样式
        barrier-line/.style={thick, red!70!black, dashed},
        % 参考线样式
        ref-line/.style={thin, gray!60, dashed}
    ]

        % 0. 绘制填充阴影 (放在最底层)
        \fill[blue!5] (0,0) 
            -- (0.2, 0.6) -- (0.4, -0.3) -- (0.6, 0.2) -- (0.8, 1.1) 
            -- (1.1, 0.3) -- (1.3, 0.7) -- (1.6, -0.5) -- (1.9, 0.2) 
            -- (2.2, 1.5) -- (2.5, 0.4) -- (2.8, 0.9) -- (3.1, -0.45) 
            -- (3.4, 0.1) -- (3.7, -0.4) -- (4.0, 0.6) -- (4.3, -0.2) -- (4.5, 0.3) % 蓝色路径结束
            -- (4.5, -0.5) -- (2.2, -0.5) -- (2.1, -0.9) -- (0.8, -0.9) -- (0.7, -1.4) -- (0.2, -1.4) -- (0, -2.0) % 红色路径逆序
            -- cycle;

        % 1. 坐标轴
        \draw[->, thick] (0, -2.5) -- (0, 2.2) node[above] {$Y_n$};
        \draw[->, thick] (0, 0) -- (5.5, 0) node[right] {$n$};

        % 2. 终点限制线 y = Delta/c (假设 Delta/c = 0.8)
        \draw[thick, black!80] (5, 0.8) -- (5.4, 0.8) node[right] {$\frac{\Delta}{c}$};
        \draw[ref-line] (0, 0.8) -- (5, 0.8);

        % 3. 时间 L 的垂线
        \draw[dashed, black!80] (4.5, -2.5) -- (4.5, 2.0) node[above] {$L \ge \Delta^{2+c}$};

        % 4. 密集随机游走路径 (零均值特征)
        % 定义一组波动点
        \draw[rw-path] (0,0) 
            -- (0.2, 0.6) -- (0.4, -0.3) -- (0.6, 0.2) -- (0.8, 1.1) 
            -- (1.1, 0.3) -- (1.3, 0.7) -- (1.6, -0.5) -- (1.9, 0.2) 
            -- (2.2, 1.5) -- (2.5, 0.4) -- (2.8, 0.9) -- (3.1, -0.45) 
            -- (3.4, 0.1) -- (3.7, -0.4) -- (4.0, 0.6) -- (4.3, -0.2) -- (4.5, 0.3);

        % 5. 动态屏障 Max Y_i - Delta (假设 Delta = 2.0)
        % 初始 Max=0 -> 屏障 -2.0
        % 第一高峰 0.6 -> 屏障 -1.4
        % 第二高峰 1.1 -> 屏障 -0.9
        % 第三高峰 1.5 (最高) -> 屏障 -0.5
        \draw[barrier-line] 
            (0, -2.0) -- (0.2, -1.4) -- (0.7, -1.4) -- (0.8, -0.9) 
            -- (2.1, -0.9) -- (2.2, -0.5) -- (4.5, -0.5) 
            node[right] {$\max_{i\le n} Y_i - \Delta$};

        % 6. 标注 Delta 间距 (在最高点处)
        \draw[ref-line] (0, 1.5) -- (2.2, 1.5);
        \draw[<->] (2.3, 1.5) -- (2.3, -0.5) node[midway, left] {$\Delta$};

        % 7. 关键点标注
        \node[below left] at (0,0) {$0$};
        \fill[blue!70!black] (4.5, 0.3) circle (1.5pt) node[right, black] {$Y_L$};

    \end{tikzpicture}}
\end{frame}

\begin{frame}{Spine-walk lemma group III: tilted tubes (Lemmas 3.6--3.7)}
    % Purpose: gives tube bounds for tilted boundaries; used in second moment bounds and Lemma 4.6.
    \small
    % Lemma `lemma:4.2` statement header (too long to fully restate; keep the role + key parameters).
\begin{block}{Lemma \texttt{lemma:4.2} (small deviation with small drift)}
Controls tube probabilities for the spine walk under a \emph{tilted} (drifting) tube with boundaries $(r_j,s_j)$; provides the key exponential upper bounds used in second-moment arguments and in \texttt{lemma:4.6}.
\end{block}


    \vspace{1mm}
    % Lemma `lemma:4.3` role + key conclusion (lower bound under tube end constraint).
\begin{block}{Lemma 3.7 (lower bound under tube end constraint)}
Provides a matching \emph{lower bound} for the expected number of particles (or tube-following paths) when additionally conditioning the tube endpoint, which is the first-moment input in Paley--Zygmund for Lemma 3.9.
\end{block}

    \vspace{1mm}
    \centering
    \resizebox{!}{0.33\textheight}{% Source: contents/part-2-eng.tex (Figure: tube constraints illustration)
\begin{tikzpicture}[
    xscale=1.2,
    yscale=1.2,
    >=Stealth,
    rw-path/.style={thick, blue!70!black},
    running-max/.style={thin, gray!80, dashed},
    barrier-red/.style={thick, red!70!black, dashed},
    barrier-green/.style={thick, green!60!black, dashed},
    ref-line/.style={thin, black!80, dashed}
]

    % shaded region
    \fill[blue!5]
        (0,0) -- (0.5, 1.1) -- (1.5, -0.8) -- (2.5, 1.6) -- (3.5, 0.4) -- (4, 1.5)
        -- (5, -0.5) -- (6, 1.7) -- (7, -0.2) -- (8, 0.5)
        -- (8, -0.5) -- (6, -0.5) -- (5.95, -0.6)-- (2.5, -0.6) -- (2.3, -1.1) -- (0.5, -1.1) -- (0, -2.2)
        -- cycle;

    % axis
    \draw[->, thick] (-0.2, -2.5) -- (-0.2, 2.5) node[above] {$Y_n$};

    % level lines
    \draw[barrier-green] (-0.3, 1.8) -- (8.2, 1.8) node[right] {$M_j$};
    \draw[black!80, dashed] (-0.3, 1.2) -- (8.2, 1.2) node[right] {$M_j - \lceil\eta\Delta\rceil$};

    % time partition lines
    \draw[ref-line] (0, -2.6) -- (0, 2.2) node[below=5.6cm] {$\ell_j$};
    \draw[ref-line] (4, -2.6) -- (4, 2.2) node[below=5.6cm] {$\ell_{j+1}$};
    \draw[ref-line] (8, -2.6) -- (8, 2.2) node[below=5.6cm] {$\ell_{j+2}$};
    \node at (2, -2.8) {$I_j$};
    \node at (6, -2.8) {$I_{j+1}$};

    % static green line for interval constraint
    \draw[barrier-green] (4, -1.0) -- (8.2, -1.0) node[right] {$(M_j - \lceil\eta\Delta\rceil) - \Delta$};

    % random walk path
    \draw[rw-path] (0,0)
        -- (0.5, 1.1) -- (1.5, -0.8) -- (2.5, 1.6) -- (3.5, 0.4) -- (4, 1.5)
        -- (5, -0.5) -- (6, 1.7) -- (7, -0.2) -- (8, 0.5);

    % running max
    \draw[running-max] (0,0) -- (0.5, 1.1) -- (2.3, 1.1) -- (2.5, 1.6) -- (5.95, 1.6) -- (6, 1.7) -- (8, 1.7);

    % dynamic red barrier
    \draw[barrier-red] (0, -2.2) -- (0.5, -1.1) -- (2.3, -1.1) -- (2.5, -0.6) -- (5.95, -0.6) -- (6, -0.5) -- (8, -0.5)
        node[right=0.25cm] {$\max_{i\le n} Y_i - \Delta$};

    % delta bracket
    \draw[<->] (6.2, 1.7) -- (6.2, -0.5);
    \node[right] at (6.2, 0.6) {$\Delta$};
\end{tikzpicture}
}
\end{frame}

\begin{frame}{How spine-walk lemmas use Mogul'skii (as in the thesis)}
    % Purpose: explain how Mogulskii is invoked inside Part 2 lemmas, before stating Mogulskii itself.
    \small
    \begin{itemize}
        \item Many-to-one reduces BRW tube counts to probabilities for the spine walk $Y$.
        \item Mogul'skii provides the sharp constant $-\pi^2/2$ for small deviations in tubes (variance $\sigma^2$).
        \item The thesis plugs Mogul'skii into Lemma 3.4 (horizontal tube), then extends to Lemma 3.5 and Lemma 3.6--3.7 by truncation / endpoint constraints.
    \end{itemize}
\end{frame}

\begin{frame}{Theorem 2.4 (Mogul'skii): statement}
    \small
    % Extracted verbatim (English part) from: contents/known-results-eng.tex (Theorem \label{theorem:mogulskii})
% Purpose: provide the sharp small-deviation exponent (and the constant -pi^2/2) for tube events.
\begin{theorem}[Mogul'skii's Theorem]\label{theorem:mogulskii-slide}
Suppose that $\mathbb E\xi_i=0$, $\mathbb E\xi_i^2=1$ and that the sequence $(x(n))$ is such that
\begin{equation*}
    \lim_{n\to\infty}x(n)=\infty,\quad
    \lim_{n\to\infty}x(n)n^{-1/2}=0.
\end{equation*}
Then the relation
\begin{equation*}
    \log\mathbb P(s_n\in G)\sim-\frac{\pi^2}{2}H_2(G)nx^{-2}(n)
\end{equation*}
holds for set $G\in\mathfrak N$.
\end{theorem}

\end{frame}

\begin{frame}{Theorem 2.4 (Mogul'skii): meaning for tube events}
    \small
    \begin{itemize}
        \item Regime: tube width $x(n)\to\infty$ but $x(n)\ll \sqrt n$ (moderate deviations).
        \item Conclusion: $\log\P(\text{stay in tube})$ is of order $-n/x(n)^2$ with explicit constant $-\pi^2/2$.
        \item In this thesis: take $x(n)\asymp \Delta$ and $n\asymp L$, so tube probabilities behave like
        $\exp(-(\pi^2/2)\sigma^2\,L/\Delta^2)$ up to sub-exponential factors.
    \end{itemize}
\end{frame}

\begin{frame}{Coupling lemma group: killed BRW vs $N$-BRW (Lemma 3.17)}
    % Purpose: connects BRW tube estimates to bounds for the $N$-BRW minimum.
    \small
    % Lemma 3.17 role in the proof of Theorem 1.1.
\begin{lemma}[Lemma 3.17]
If the killing rule is adapted and $\mathcal V^{\mathrm k}_0\subseteq\mathcal V^N_0$, then one can couple $(\mathcal V^N,\mathcal V^{\mathrm k})$ so that $\mathcal M^N_k\le\mathcal M^{\mathrm k}_k$ for all $k\le\tau:=\inf\{k:\#\mathcal Z^{\mathrm k}_k>N\}$.
\end{lemma}

\end{frame}

\begin{frame}{BRW second-moment lemma chain: Assumption 6 propagation}
    % Purpose: emphasize how Assumption 6 is transmitted through the lemmas.
    \small
    Assumption 6:
    \[
        \log \E\!\left[
            \sum_{i,j\in\N,\ i\neq j}
            e^{-(\Xi_i+\Xi_j)}
        \right]
        <\infty.
    \]
    \vspace{-1mm}
    \begin{itemize}
        \item Lemma 3.8: \textbf{uses Assumption 6} to bound the cross term $\sum_{i\neq j}\Sigma(z^i)\Sigma(z^j)$.
        \item Lemma 3.9: applies \textbf{Paley--Zygmund} and \textbf{invokes Lemma 3.8} (hence Assumption 6).
        \item Lemma 3.10: invokes \textbf{Lemma 3.9} (hence Lemma 3.8 + Assumption 6).
    \end{itemize}
\end{frame}

\begin{frame}{BRW lemma group: Lemma 3.8 (second moment bound; uses Assumption 6)}
    % Purpose: key 2nd moment estimate; keep the formula.
    \footnotesize
    % Lemma `lemma:4.5` statement (verbatim math, English-only wrapper).
\begin{lemma}[\texttt{lemma:4.5}]
Let $v>0$, $0<d<\frac{\pi^2}2\sigma^2$, $c\ge 0$, and $r_j\le s_j$ (as defined in the thesis right before this lemma).
There exists $\Delta_0>0$ such that for all $\Delta\ge\Delta_0$, any $\ell\ge 1$ and $\delta>r_\ell$,
\begin{equation*}
    \mathbb E\!\left[
        \left(
            \sum_{|x|=\ell}
            {\bf 1}_{
                \{
                    r_j\le\mathcal V(x_j)\le s_j,\,
                    \forall j\in[\![1,\ell]\!],\,
                    \mathcal V(x)\le\delta
                \}
            }
        \right)^2
    \right]
    \le
    e^{4\Delta^v+2\delta-2d\frac{\ell-1}{\Delta^2}}
    \sum_{m=0}^{\ell-1}
    e^{d\frac m{\Delta^2}-r_m}
    +
    e^{\delta+\Delta^v-\frac d{\Delta^2}\ell}.
\end{equation*}
\end{lemma}

\end{frame}

\begin{frame}{BRW lemma group: Lemma 3.8 (proof idea; Assumption 6)}
    % Purpose: outline the proof approach exactly as in contents/part-3-eng.tex.
    \small
    \begin{itemize}
        \item Let $F_x$ be the tube indicator at generation $\ell$.
        \item Expand $\left(\sum_{|x|=\ell}F_x\right)^2$ and group pairs $(x,y)$ by their most recent common ancestor $z$ at generation $m$.
        \item The cross term is a sum over $i\neq j$ of $\Sigma(z^i)\Sigma(z^j)$ (two different children): Assumption 6 controls its expectation.
    \end{itemize}
\end{frame}

\begin{frame}{BRW lemma group: Lemma 3.9 (tube existence; uses Lemma 3.8)}
    % Purpose: tube existence at $L^{1/3}$ scale.
    \footnotesize
    % Lemma `lemma:4.6` statement (verbatim math, English-only wrapper).
\begin{lemma}[Lemma 3.9]
Let $c>0$. For $L\ge 1$, set $r_j:=\frac{\pi^2\sigma^2}{2c^2}L^{-2/3}j$ and $s_j:=r_j+cL^{1/3}$ for $j\ge 0$.
For any $v>0$,
\[
    \lim_{L\to\infty}
    L^{-\frac 13}
    \log\mathbb P\left(
        \exists |x|=L:
        \begin{aligned}
        r_j\le\mathcal V(x_j)\le s_j,\,
        \forall j\in[\![1,L]\!],\\
        \mathcal V(x)\le r_L+vL^{\frac 13}
        \end{aligned}
    \right)
    =
    0,
\]
and moreover,
\[
    \lim_{L\to\infty}
    L^{-\frac 13}
    \log\mathbb E\!\left[
        \sum_{1\le |x|\le L}
        {\bf 1}_{
            \{
                r_j\le\mathcal V(x_j)\le s_j,\,
                \forall j\in[\![1,|x|]\!]
            \}
        }
    \right]
    =
    c.
\]
\end{lemma}

\end{frame}

\begin{frame}{BRW lemma group: Lemma 3.9 (proof organization)}
    % Purpose: keep the original lemma dependencies.
    \small
    \begin{itemize}
        \item Prove the lower bound via \textbf{Paley--Zygmund}.
        \item First moment lower bound: Lemma 3.7.
        \item Second moment upper bound: \textbf{Lemma 3.8} (uses Assumption 6) + Lemma 3.6.
        \item Conclude the comparison by combining Lemma 3.6 (upper) and Lemma 3.7 (lower).
    \end{itemize}
\end{frame}

\begin{frame}{BRW lemma group: Lemma 3.10 (log-scale tube bounds; uses Lemma 3.9)}
    % Purpose: log-scale tube probability / expectation asymptotics.
    \footnotesize
    % Lemma `lemma:4.7` statement (verbatim math, English-only wrapper).
\begin{lemma}[\texttt{lemma:4.7}]
Let $b>v>0$ and $r:=-b\log n$. Then
\[
    \lim_{n\to\infty}
    \frac{1}{\log n}
    \log\mathbb P\!\left(
        \exists |x|=\lfloor\log^{5/2}n\rfloor:
        \begin{aligned}
            r\le\mathcal V(x_j)\le 0,\,
            \forall j\in[\![1,|x|]\!],\\
            \mathcal V(x)\le r+v\log n
        \end{aligned}
    \right)
    =
    v-b,
\]
moreover,
\[
    \mathbb E\!\left[
        \sum_{\ell=1}^\infty
        \sum_{|x|=\ell}
        {\bf 1}_{
            \{
                r\le\mathcal V(x_j)\le 0,\,
                \forall j\in[\![1,\ell]\!]
            \}
        }
    \right]
    =
    O_n(\log n),
\]
and
\[
    \mathbb E\!\left[
        \sum_{|x|=\lfloor\log^{5/2}n\rfloor}
        {\bf 1}_{
            \{
                r\le\mathcal V(x_j)\le 0,\,
                \forall j\in[\![1,|x|]\!],\,
                \mathcal V(x)\le r+v\log n
            \}
        }
    \right]
    =
    n^{-b+v+o_n(1)}.
\]
\end{lemma}


\end{frame}

\begin{frame}{BRW lemma group: Lemma 3.10 (proof organization)}
    % Purpose: keep the original dependencies and proof flow.
    \small
    \begin{itemize}
        \item Upper estimate: combine Lemma 3.4 (upper) and Lemma 3.7 (lower) with a tighter boundary.
        \item Lower estimate: apply Paley--Zygmund and invoke \textbf{Lemma 3.9} (hence Lemma 3.8 + Assumption 6).
        \item The remaining bound follows from standard random-walk arguments.
    \end{itemize}
\end{frame}

\begin{frame}{Other values of $a$: $a=1$ and $a=0$}
    % Purpose: address how the proof extends to endpoints, per the thesis.
    \small
    \begin{itemize}
        \item $a=1$: obtained by the same scheme (tail bounds + 4th moment + Cauchy--Schwarz), using the already proved $a=1$ case in Theorem 1.1 when applying the speed proof.
        \item $a=0$: the statement of Theorem 1.1 additionally requires the exponential moment condition on $\Xi_1$ (Assumption 9).
    \end{itemize}
\end{frame}

\begin{frame}{Left tail lower bound (final step): uses no-large-drop small deviation}
    % Purpose: conclude with the left tail bound and the required lemma + figure.
    \small
    The left tail bound
    \[
        \P(\mathcal M^N_{\ell_N}\le \zeta\log N)\le N^{-\varepsilon}
    \]
    is proved for all $a\in[0,1]$ by reducing to a no-large-drop event for the spine walk and applying Lemma 3.5.
    \vspace{1mm}
    \centering
    \resizebox{0.78\textwidth}{!}{\begin{tikzpicture}[
        xscale=1.7, % 修改这里控制横向（代数 n）的间距
        yscale=1.1, % 修改这里控制纵向（空间位置 V）的伸展程度
        >=Stealth,
        % 密集随机游走样式
        rw-path/.style={thick, blue!70!black},
        % 动态屏障样式
        barrier-line/.style={thick, red!70!black, dashed},
        % 参考线样式
        ref-line/.style={thin, gray!60, dashed}
    ]

        % 0. 绘制填充阴影 (放在最底层)
        \fill[blue!5] (0,0) 
            -- (0.2, 0.6) -- (0.4, -0.3) -- (0.6, 0.2) -- (0.8, 1.1) 
            -- (1.1, 0.3) -- (1.3, 0.7) -- (1.6, -0.5) -- (1.9, 0.2) 
            -- (2.2, 1.5) -- (2.5, 0.4) -- (2.8, 0.9) -- (3.1, -0.45) 
            -- (3.4, 0.1) -- (3.7, -0.4) -- (4.0, 0.6) -- (4.3, -0.2) -- (4.5, 0.3) % 蓝色路径结束
            -- (4.5, -0.5) -- (2.2, -0.5) -- (2.1, -0.9) -- (0.8, -0.9) -- (0.7, -1.4) -- (0.2, -1.4) -- (0, -2.0) % 红色路径逆序
            -- cycle;

        % 1. 坐标轴
        \draw[->, thick] (0, -2.5) -- (0, 2.2) node[above] {$Y_n$};
        \draw[->, thick] (0, 0) -- (5.5, 0) node[right] {$n$};

        % 2. 终点限制线 y = Delta/c (假设 Delta/c = 0.8)
        \draw[thick, black!80] (5, 0.8) -- (5.4, 0.8) node[right] {$\frac{\Delta}{c}$};
        \draw[ref-line] (0, 0.8) -- (5, 0.8);

        % 3. 时间 L 的垂线
        \draw[dashed, black!80] (4.5, -2.5) -- (4.5, 2.0) node[above] {$L \ge \Delta^{2+c}$};

        % 4. 密集随机游走路径 (零均值特征)
        % 定义一组波动点
        \draw[rw-path] (0,0) 
            -- (0.2, 0.6) -- (0.4, -0.3) -- (0.6, 0.2) -- (0.8, 1.1) 
            -- (1.1, 0.3) -- (1.3, 0.7) -- (1.6, -0.5) -- (1.9, 0.2) 
            -- (2.2, 1.5) -- (2.5, 0.4) -- (2.8, 0.9) -- (3.1, -0.45) 
            -- (3.4, 0.1) -- (3.7, -0.4) -- (4.0, 0.6) -- (4.3, -0.2) -- (4.5, 0.3);

        % 5. 动态屏障 Max Y_i - Delta (假设 Delta = 2.0)
        % 初始 Max=0 -> 屏障 -2.0
        % 第一高峰 0.6 -> 屏障 -1.4
        % 第二高峰 1.1 -> 屏障 -0.9
        % 第三高峰 1.5 (最高) -> 屏障 -0.5
        \draw[barrier-line] 
            (0, -2.0) -- (0.2, -1.4) -- (0.7, -1.4) -- (0.8, -0.9) 
            -- (2.1, -0.9) -- (2.2, -0.5) -- (4.5, -0.5) 
            node[right] {$\max_{i\le n} Y_i - \Delta$};

        % 6. 标注 Delta 间距 (在最高点处)
        \draw[ref-line] (0, 1.5) -- (2.2, 1.5);
        \draw[<->] (2.3, 1.5) -- (2.3, -0.5) node[midway, left] {$\Delta$};

        % 7. 关键点标注
        \node[below left] at (0,0) {$0$};
        \fill[blue!70!black] (4.5, 0.3) circle (1.5pt) node[right, black] {$Y_L$};

    \end{tikzpicture}}
\end{frame}

% --- Closing slide ---
\begin{frame}
    \begin{center}
        \Huge \textcolor{blue!80!black}{Thank you!}

        \vspace{1cm}
        \Large Questions and comments are welcome.
    \end{center}
\end{frame}

\end{document}
